% ===========================================================================
% Chapter 2: Randomness
% ===========================================================================
% Gacs source lines 2325-3756 of
% Gacs_2021_Descriptional_Complexity_and_Randomness.tex, reproduced verbatim with
% the author's permission.  Edits to his text: sloppypar wrappers dropped, runs of
% blank lines collapsed, the proofs of Theorems 2.3.3, 2.3.4, 2.3.5 and Corollary
% 2.3.30 given a title or an empty body so that the dependency graph can attach
% them, and references to his Appendix A replaced by their printed numbers.
% Section 4.1 (sec:cont-spaces) is referenced by label.
%
% The Lean formalization is Gacs_Chapter_2/Randomness/ (namespace
% Kolmogorov.Randomness); Gacs_Chapter_2/strategy.md has the plan, the design
% decisions and the progress log.
%
% Numbering: within each section his theorems run on one counter and everything
% else on another (macros/common.tex), so every number below is the one printed in
% his manuscript.  Numbered items he did not label carry labels <kind>:g<number>,
% e.g. def:g2.3.15.  Statements that are ours use ourdefinition/ourtheorem (labels
% def:ch2-..., thm:ch2-...), which have their own counter.

\chapter{Randomness}

Chapter 2 of G\'acs' manuscript defines randomness: for finite strings under the uniform
distribution (Section~\ref{sec:tests}) and under a computable distribution
(Section~\ref{sec:discrand}), and for infinite sequences. All of it is formalized in Lean,
with two exceptions. Theorem~\ref{thm:g2.3.6} on monotone machines is planned. The general
measure theory of Subsections 2.3.2--2.3.4 is not formalized, because Mathlib provides it
and no result of the chapter needs it. His text is reproduced with his permission; the
annotations, the dependency edges and the formalization notes are ours.

Three conventions apply throughout.
\begin{itemize}
\item A test \( d \) is carried as \( t = 2^{d} \), with values in \( [0,\infty] \).
  Universality \( d \lea d_{0} \) becomes \( t \le 2^{c} t_{0} \), and the deficiency
  \( -\log P(x) - K(x) \) becomes \( 2^{-K(x)}/P(x) \). No logarithm appears in the Lean
  statements.
\item The constants \( K(d) + K(P) \) of Theorems~\ref{thm:univ-integrable-test},
  \ref{thm:simple-rand-cons-det}, \ref{thm:g2.2.3} and~\ref{thm:univ-integrable-test-seq}
  are formalized twice: with a constant that depends on the test and the measure, and with
  his explicit constant, where \( K(d) \) is read through a description of \( d \)
  (Definition~\ref{def:ch2-description}).
\item Sequences are binary. Every statement about them is made through finite strings:
  cylinders, open sets given by generating strings, measures by their values on cylinders,
  tests by their values on prefixes.
\end{itemize}
As in Chapter~1, \( K \) is the prefix complexity of an optimal prefix interpreter \( U \)
that every result takes as a hypothesis (\texttt{IsOptimalPrefixConditional U}).

\section{Uniform distribution} \label{sec:tests}

Speaking of a ``random binary string'', one generally understands
randomness with respect to the coin-toss distribution (when the
probability $P(x)$ of a binary sequence $x$ of length $n$ is $2^{-n}$).
This is the only distribution considered in the present section.

One can hope to distinguish sharply between random and nonrandom
\emph{infinite sequences}.
Indeed, an infinite binary sequence whose
elements are 0 with only finitely many exceptions, can be considered
nonrandom without qualification.
This section defines randomness for finite strings.
For a sequence $x$ of length $n$, it would be
unnatural to fix some number $k$, declare $x$ nonrandom when its
first $k$ elements are 0, but permit it to be random if only the
first $k-1$ elements are 0.
So, we will just declare some finite strings less random than others.
For
this, we introduce a certain real-valued function $d(x) \ge 0$ measuring
the \df{deficiency of randomness} in the string $x$.
The occurrence
of a nonrandom sequence is considered an  exceptional event,
therefore the function $d(x)$ can assume large values only with small
probability.
What is ``large'' and ``small'' depends only on our ``unit of measurement''.
We require for all $n,k$
 \begin{equation} \label{eq:testcond}
\sum\setof{ P(x): x\in \dB^{n},\ d(x) > k } < 2^{-k},
 \end{equation}
saying that there be at most $2^{n-k}$ binary sequences $x$ of length
$n$ with $d(x)>k$.
Under this condition, we even allow $d(x)$ to take the value $\infty$.

To avoid arbitrariness in the distinction between random and
nonrandom, the function $d(x)$ must be \emph{simple}.
We assume
therefore that the set $\setof{ (n,k,x) : \len{x} = n,\; d(x)>k }$ is
recursively enumerable, or, which is the same, that the function
 \begin{align*}
  \d : \Og\to\rint{-\infty}{\infty}
 \end{align*}
is lower semicomputable.

% Gacs source line 2370
 \begin{remark}
We do not assume that the set of strings $x \in \dB^{n}$ with \emph{small}
deficiency of randomness is enumerable, because then it would be easy
to construct a ``random'' sequence.
 \end{remark}

% Gacs source line 2376
\begin{definition}
  \label{def:g2.1.2}
  \lean{Kolmogorov.Randomness.IsMLTestUniform, Kolmogorov.Randomness.IsUniversalMLTestUniform}
  \leanok
  \uses{def:semicomputability}
A lower semicomputable function $d : \dS_{2} \to \dR$ (where $\dR$ is the
set of real numbers) is called a \df{Martin-L\"of test (ML-test)}, or a
\df{probability-bounded test} if it satisfies~\eqref{eq:testcond}.

A ML-test $d_{0}(x)$ is \df{universal} if it additively dominates all
other ML-tests: for any other ML-test $d(x)$ there is a
constant $c<\infty$ such that for all $x$ we have
$d(x) < d_{0}(x) + c$.
\end{definition}

\begin{formalization}
A test is carried as \( t = 2^{d} \), a lower semicomputable function from strings to
\( [0,\infty] \). Condition~\eqref{eq:testcond} reads
\( \card{\setof{x \in \dB^{n} : 2^{k} < t(x)}} \cdot 2^{k} \le 2^{n} \) for all
\( n, k \in \dN \). It is stated with \( \le \): his text says ``at most \( 2^{n-k} \)'', and
that is what the proof of Theorem~\ref{thm:MLcompl} uses. Universality is
\( t \le 2^{c} t_{0} \) with \( c \in \dN \).
\end{formalization}

If a test $d_{0}$ is universal then any other test $d$ of randomness
can discover at most by a constant amount more deficiency of
randomness in any sequence $x$ than $d_{0}(x)$.
Obviously, the
difference of any two universal ML-tests is bounded by some
constant.

The following theorem reveals a simple connnection between descriptional
complexity and a certain randomness property.

% Gacs source line 2397
 \begin{definition}
  \label{def:g2.1.3}
  \lean{Kolmogorov.Randomness.mlDeficiency}
  \leanok
  \uses{def:fixed-U}
Let us define the following function for binary strings $x$:
 \begin{align}\label{eq:d0}
 d_{0}(x) = \len{x} - C(x \mvert\; \len{x}).
 \end{align}
 \end{definition}

\begin{formalization}
In multiplicative form: \texttt{mlDeficiency U x} \( = 2^{\len{x}} \cdot 2^{-C(x \mvert \len{x})} \),
with the length passed to \( C \) as its binary representation.
\end{formalization}

% Gacs source line 2404
\begin{theorem}[Martin-L\"of] \label{thm:MLcompl}
  \lean{Kolmogorov.Randomness.mlDeficiency_isUniversalMLTestUniform, Kolmogorov.Randomness.mlDeficiency_isMLTestUniform, Kolmogorov.Randomness.mlTest_le_mlDeficiency}
  \leanok
  \uses{def:g2.1.2, def:g2.1.3}
 The function $d_{0}(x)$ is a universal Martin-L\"of test.
 \end{theorem}
\begin{proof}
  \proves{thm:MLcompl}
  \leanok
  \uses{thm:Ksemicp, thm:Kmajor}
 Since $C(x \mvert y)$ is semicomputable from above, $d_{0}$ is
semicomputable from below.
The property of semicomputability holds
for $d_{0}$ as a straightforward consequence of Theorem~\ref{thm:Ksemicp}
Therefore $d_{0}$ is a ML-test.
We must show that it is larger (to
within an additive constant) than any other ML-test. Let $d$ be a
ML-test.
Let us define the function $F(x,y)$ to be $y - d(x)$ for
$y = \len{x}$, and $\infty$ otherwise.
Then $F$ is upper semicomputable, and
satisfies~\eqref{eq:Kmajor}.
Therefore by Theorem~\ref{thm:Kmajor}, we
have $C(x \mvert \len{x}) \lea \len{x} - d(x)$.
 \end{proof}

\begin{formalization}
The proof is his, through Levin's Theorem~\ref{thm:Kmajor}. In multiplicative form
\( F(x, \len{x}) \) is the least \( m \) with \( 2^{\len{x}} < 2^{m} t(x) \). Its minimality
gives \( 2^{F(x,\len{x})} t(x) \le 2^{\len{x}+1} \), so the constant of the theorem is one more
than the constant of Levin's theorem.
\end{formalization}

Theorem~\ref{thm:MLcompl} says that under very general assumptions
about randomness, those strings $x$ are random whose descriptional
complexity is close to its maximum, $\len{x}$.
The more ``regularities'' are
discoverable in a sequence, the less random it is.
However, this is
true only of regularities which decrease the descriptional
complexity.
``Laws of randomness'' are regularities whose probability
is high, for example the law of large numbers, the law of iterated
logarithm, the arcsine law, etc.
A random sequence will satisfy \emph{all such laws}.

% Gacs source line 2437
 \begin{example}
The law of large numbers
says that in most binary sequences of length $n$, the number of 0's
is close to the numer of 1's.
We prove this law of probability theory
by constructing a ML-test $d(x)$ taking large values on sequences in
which the number of 0's is far from the number of 1's.

Instead of requiring~\eqref{eq:testcond} we require a somewhat stronger
property of $d(x)$: for all $n$,
 \begin{equation}  \label{eq:llntestcd}
 \sum_{x\in \dB^{n}} P(x)2^{d(x)} \le 1.
 \end{equation}

From this, the inequality~\eqref{eq:testcond} follows by the following
well-known inequality called the Markov Inequality which says the
following.
Let $P$ be any probability distribution, $f$ any
nonnegative function, with expected value $E_{P}(f)=\sum_{x}P(x)f(x)$.
For all $\lg\ge 0$ we have
 \begin{equation} \label{eq:Markov}
 \sum\setof{ P(x): f(x) > \lg E_{P}(f)} \le \lg.
 \end{equation}

For any string $x\in \dS$, and natural number $i$, let $N(i \mvert x)$
denote the number of occurrences of $i$ in $x$.
For a binary string $x$
of length $n$ define $p_{x} = N(1 \mvert x)/n$, and
 \begin{align*}
   P_{x}(y)	&= p_{x}^{N(1 \mvert y)}(1-p_{x})^{N(0 \mvert y)}
 \\ d(x)	&= \log P_{x}(x) + n - \log(n+1).
 \end{align*}
  We show that $d(x)$ is a ML-test.
It is obviously computable.
We prove~\eqref{eq:llntestcd}.
We have
 \begin{align*}
  \sum_{x} P(x) 2^{d(x)} &= \sum_{x} 2^{-n} 2^{n} P_{x}(x)\frac{1}{n+1}
	= \frac{1}{n+1}\sum_{x} P_{x}(x)
  \\ &= \frac{1}{n+1}\sum_{k=0}^{n}\binom{n}{k}(\frac{k}{n})^k
   (1-\frac{k}{n})^{n-k} < \frac{1}{n+1}\sum_{k=0}^{n} 1 = 1.
 \end{align*}
The test $d(x)$ expresses the (weak) law of large numbers in a rather
strong version. We rewrite $d(x)$ as
 \[
 d(x) = n(1 - h(p_{x})) - \log(n+1)
 \]
 where $h(p) = -p\log p - (1-p)\log(1-p)$.
The entropy function
$h(p)$ achieves its maximum 1 at $p=1/2$.
Therefore the test $d(x)$
tells us that the probability of sequences $x$ with $p_{x} < p < 1/2$
for some constant $p$ is bounded by the exponentially decreasing
quantitiy $(n+1)2^{-n(1-h(p))}$.
We also see that since for some
constant $c$ we have
 \[
   1-h(p) > c(p-1/2)^2,
 \]
therefore if the difference $| p_{x} - 1/2 |$ is much larger than
 \[
   \sqrt\frac{\log n}{cn}
 \]
  then the sequence $x$ is ``not random'', and hence, by
Theorem~\ref{thm:MLcompl} its complexity is significantly smaller than $n$.
 \end{example}

\begin{formalization}
Not formalized. Markov's inequality~\eqref{eq:Markov} is proved where it is used: in
Proposition~\ref{propo:ml-integrable}, Corollary~\ref{cor:g2.2.8} and
Theorem~\ref{thm:g2.3.2}.
\end{formalization}

\section{Computable distributions}
 \label{sec:discrand}

Let us generalize the notion of randomness to arbitrary discrete computable
probability distributions.

\subsection{Two kinds of test}

Let $P(x)$ be a probability distribution over some discrete countable
space $\Og$, which we will identify for simplicity with the set $\dS$ of
finite strings of natural numbers.
Thus, $P$ is a nonnegative function with
 \[
 \sum_{x} P(x) = 1.
 \]
  Later we will consider probability distributions $P$ over the space
$\dN^{\dN}$, and they will also be characterized by a nonnegative function
$P(x)$ over $\dS$.
However, the above condition will be replaced by a different one.

We assume $P$ to be computable, with some G\"odel number $e$.
We want to define a test $d(x)$ of randomness with respect to $P$.
It will
measure how justified is the assumption that $x$ is the outcome of an
experiment with distribution $P$.

\begin{ourdefinition}[Computable probability distribution]
  \label{def:ch2-computable-measure}
  \lean{Kolmogorov.Randomness.ComputableMeasure, Kolmogorov.Randomness.IsComputableENNReal}
  \leanok
  \uses{def:computable-function, def:semimeasure}
A \df{computable probability distribution} on binary strings is a function \( P \) from
strings to \( [0,\infty] \) with \( \sum_{x} P(x) = 1 \) for which there are computable
rational functions \( \ell_{k}(x) \le P(x) \le u_{k}(x) \) with
\( u_{k}(x) - \ell_{k}(x) \le 2^{-k} \).
\end{ourdefinition}
\begin{definedinlean}\proves{def:ch2-computable-measure}\leanok\end{definedinlean}

\begin{formalization}
This definition is ours. It fixes what Section~\ref{sec:discrand} assumes about \( P \):
strings are binary, and the G\'odel number \( e \) is replaced by the two approximating
sequences. Chapter~1's \texttt{Kolmogorov.IsComputableENNReal} asks for one dyadic
approximation within \( 2^{-k} \). The two-sided version makes both \( P \) and \( 1/P \)
lower semicomputable directly, and the universal test \( 2^{-K(x)}/P(x) \) needs the second.
\end{formalization}

% Gacs source line 2533
 \begin{definition}
  \label{def:g2.2.1}
  \lean{Kolmogorov.Randomness.IsIntegrableTest, Kolmogorov.Randomness.IsUniversalIntegrableTest}
  \leanok
  \uses{def:ch2-computable-measure, def:semicomputability}
A function $d:\dS\to\dR$ is an \df{integrable test}, or
\df{expectation-bounded test} of randomness
with respect to $P$ if it is lower semicomputable and satisfies the condition
 \begin{equation} 	\label{eq:newtestcond}
\sum_{x} P(x)2^{d(x)} \le 1.
 \end{equation}
It is \df{universal} if it dominates all other integrable tests to within an
additive constant.
 \end{definition}

\begin{formalization}
Multiplicatively, \texttt{IsIntegrableTest P t} says that \( t = 2^{d} \) is lower
semicomputable and \( \sum_{x} P(x) t(x) \le 1 \). Universality is \( t \le 2^{c} t_{0} \).
\end{formalization}

(A similar terminology is used in~\cite{LiViBook97} in order to distinguish
tests satisfying condition~\eqref{eq:newtestcond} from Martin-L\"of tests.)

% Gacs source line 2547
\begin{remark}
  We must allow $d(x)$ to take negative values, even if only, say values
$\ge -1$, since otherwise condition~\eqref{eq:newtestcond} could only be satisfied
with $d(x)=0$ for all $x$ with $P(x)>0$.
\end{remark}

% Gacs source line 2553
 \begin{proposition}\label{propo:ml-integrable}
  \lean{Kolmogorov.Randomness.IsMLTest, Kolmogorov.Randomness.IsIntegrableTest.isMLTest, Kolmogorov.Randomness.IsMLTest.mlCorrected_isIntegrableTest, Kolmogorov.Randomness.le_mul_mlCorrected}
  \leanok
  \uses{def:g2.2.1}
Let $c = \log(2\pi^{2}/6)$.
If a function $d(x)$ satisfies~\eqref{eq:newtestcond} then it
satisfies~\eqref{eq:testcond}.
If $d(x)$ satisfies~\eqref{eq:testcond} then $d - 2\log d - c$
satisfies~\eqref{eq:newtestcond}.
 \end{proposition}
 \begin{proof}
  \proves{propo:ml-integrable}
  \leanok
The first statement follows by Markov's Inequality~\eqref{eq:Markov}.
The second one can be checked by immediate computation.
 \end{proof}

\begin{formalization}
\texttt{IsMLTest P t} is condition~\eqref{eq:testcond} for an arbitrary \( P \):
\( P\setof{x : t(x) > 2^{k}} \le 2^{-k} \) for all \( k \in \dN \). The first statement is
Markov's inequality. For the second, the correction uses the weights
\( w_{k} = 2^{k}/((k+1)(k+2)) \), for which \( \sum_{k} w_{k} 2^{-k} \) telescopes to exactly
\( 1 \). The corrected test \( t'(x) = \sup\setof{w_{k} : 2^{k} < t(x)} \) is integrable with
no free constant, and \( t \le 2(k+1)(k+2)\, t' \) on the level \( 2^{k} < t \le 2^{k+1} \).
Additively, \( d' \ge d - 2\log(d+2) - 1 \) instead of \( d - 2\log d - c \) with
\( c = \log(2\pi^{2}/6) \).
\end{formalization}

Thus, condition~\ref{eq:newtestcond} is nearly equivalent
to condition~\eqref{eq:testcond}: to each test $d$ according to one of them,
there is a test $d'$ asymptotically equal to $d$ satisfying the other.
But condition~\eqref{eq:newtestcond} has the advantage
of being just one inequality instead of the infinitely many ones.

\subsection{Randomness via complexity}

We want to find a universal integrable test.
Let us introduce the function $w(x) = P(x)2^{d(x)}$.
The above conditions on $d(x)$
imply that $w(x)$ is semicomputable from below
and satisfies $\sum_{x} w(x) \le 1$: so, it is a constructive semimeasure.
With the universal constructive semimeasure of Theorem~\ref{thm:univsem},
we also obtain a universal test for any computable probability distribution $P$.

% Gacs source line 2581
\begin{definition}
  \label{def:g2.2.4}
  \lean{Kolmogorov.Randomness.canonicalIntegrableTest}
  \leanok
  \uses{def:m, def:K, def:ch2-computable-measure}
For an arbitrary measure $P$ over a discrete space $\Og$, let us denote
 \begin{equation}\label{eq:d_c}
  \ol\d_{P}(x) = \log\frac{\m(x)}{P(x)} = -\log P(x) - K(x).
 \end{equation}
\end{definition}

\begin{formalization}
The test is carried as \( 2^{\ol\d_{P}(x)} = 2^{-K(x)}/P(x) \) (\texttt{canonicalIntegrableTest U P}),
with \( 2^{-K(x)} \) in place of \( \m(x) \). By the coding theorem
(Theorem~\ref{thm:coding}) the two differ by a bounded factor. It is defined for computable
\( P \) only, and its value is \( \infty \) where \( P(x) = 0 \).
\end{formalization}

The constants \( K(d) \) and \( K(P) \) in the next theorem need a formal reading. The
following definition and theorem are ours and supply it.

\begin{ourdefinition}[Description of a lower semicomputable function]
  \label{def:ch2-description}
  \lean{Kolmogorov.Randomness.IsLSCDescription, Kolmogorov.Randomness.descVal, Kolmogorov.Randomness.exists_lscDescription}
  \leanok
  \uses{def:semicomputability}
A string \( y \) \df{describes} a lower semicomputable function \( f \) from strings to
\( [0,\infty] \) if the approximation with index \( y \), in a fixed effective enumeration of
all lower semicomputable approximations, converges to \( f \). Every lower semicomputable
function has a description. For a test \( d \), a distribution \( P \), a transition \( T \)
and a computable map \( f \), the quantities \( K(d) \), \( K(P) \), \( K(T) \), \( K(f) \) are
read as \( K(y) \) for a description \( y \) of \( 2^{d} \), of \( P \), of
\( (x,y) \mapsto T(x,y) \) on encoded pairs, and of the graph indicator
\( (x,y) \mapsto 1[y = f(x)] \), respectively.
\end{ourdefinition}
\begin{definedinlean}\proves{def:ch2-description}\leanok\end{definedinlean}

\begin{formalization}
The enumeration is the repository's \texttt{approxEnum}, made monotone (\texttt{enumApprox}),
with \( y \) read as the number \texttt{bitsToNat y}. Every statement with an explicit
constant is quantified over all descriptions, so it holds for a shortest one, which is how
\( K(d) \) is meant. G\'acs reads \( K(d) \) as the length of a shortest program for \( d \);
the two readings should agree up to an additive constant, but that comparison is not
formalized.
\end{formalization}

\begin{ourtheorem}[Exact domination]
  \label{thm:ch2-exact-domination}
  \lean{Kolmogorov.Randomness.indexed_domination, Kolmogorov.Randomness.indexed_domination_trim, Kolmogorov.Randomness.semimeasure_le_exact, Kolmogorov.Randomness.exists_pairCode_weight, Kolmogorov.Randomness.contIndexed_domination}
  \leanok
  \uses{def:ch2-description, def:semimeasure, def:K}
Let \( G(x,z) \) be lower semicomputable in both arguments. There is a constant \( c \) such
that for every \( z \) for which \( G(\cdot,z) \) is a semimeasure, and for all \( x \),
 \[
   2^{-K(z)} G(x,z) \le 2^{c} \cdot 2^{-K(x)}.
 \]
In particular \( 2^{-K(y)} \nu(x) \le 2^{c} \cdot 2^{-K(x)} \) for every constructive
semimeasure \( \nu \) and every description \( y \) of \( \nu \), with \( c \) independent of
\( \nu \). Moreover \( 2^{-K(y)} 2^{-K(y')} \le 2^{c} \cdot 2^{-K(\ang{y,y'})} \) for all
\( y, y' \).
\end{ourtheorem}

\begin{proof}
  \proves{thm:ch2-exact-domination}
  \leanok
  \uses{thm:univsem, thm:coding, l:kraft}
Trim each \( G(\cdot,z) \), uniformly in \( z \), into a semimeasure that equals
\( G(\cdot,z) \) whenever that is already one, as in the proof of
Theorem~\ref{thm:univsem}. The mixture \( \sum_{z} 2^{-K(z)} G(x,z) \) is then a constructive
semimeasure by Kraft's inequality, and the coding theorem bounds it by
\( 2^{c} \cdot 2^{-K(x)} \).
\end{proof}

\begin{formalization}
The second statement is G\'acs' exact domination~\eqref{eq:exact-domin}
(Theorem~\ref{thm:exact-domin}), with \( 2^{-K(y)} \) in place of \( \m(\nu) \) and
\( 2^{-K(x)} \) in place of \( \m(x) \). The third is \( K(\ang{y,y'}) \lea K(y) + K(y') \).
Together they give every domination argument of this chapter the constant
\( K(d) + K(P) \): index the semimeasure by \( z = \ang{y_{d}, y_{P}} \). The continuous
analogue, against the universal semimeasure \( M \) of Theorem~\ref{thm:g2.3.5}, is
\texttt{contIndexed\_domination}.
\end{formalization}

The following theorem shows that randomness can be tested by checking how close
is the complexity $K(x)$ to its upper bound $-\log P(x)$.

% Gacs source line 2591
\begin{theorem}\label{thm:univ-integrable-test}
  \lean{Kolmogorov.Randomness.canonicalIntegrableTest_isUniversal, Kolmogorov.Randomness.canonicalIntegrableTest_isIntegrableTest, Kolmogorov.Randomness.integrableTest_le_canonical, Kolmogorov.Randomness.universalSemimeasure_div_isUniversal, Kolmogorov.Randomness.integrableTest_le_canonical_exact, Kolmogorov.Randomness.integrableTest_le_canonical_exact'}
  \leanok
  \uses{def:g2.2.1, def:g2.2.4, def:ch2-description}
The function $\ol\d_{P}(x)=-\log P(x)-K(x)$
is a universal integrable test for any fixed computable probability
distribution $P$.
More exactly, it is lower
semicomputable, satisfies~\eqref{eq:newtestcond} and
for all integrable tests $d(x)$ for $P$, we have
 \begin{equation*}d(x) \lea \ol\d_{P}(x) + K(d) + K(P).
 \end{equation*}
 \end{theorem}
 \begin{proof}
  \proves{thm:univ-integrable-test}
  \leanok
  \uses{thm:ch2-exact-domination, thm:coding, l:kraft}
Let $d(x)$ be an integrable test for $P$.
Then $\nu(x) = 2^{d(x)}P(x)$ is a constructive semimeasure, and it
has a self-delimiting program of length $\lea K(d) + K(P)$.
It follows (using the definition~\eqref{eq:m(nu)})
that $m(\nu) \lem 2^{K(d) + K(P)}$; hence inequality~\eqref{eq:exact-domin}
gives
 \[
   \nu(x) 2^{K(d) + K(P)} \lem \m(x).
 \]
Taking logarithms finishes the proof.
 \end{proof}

\begin{formalization}
Formalized in two forms.
\begin{itemize}
\item With a constant that depends on \( d \) and \( P \): \( t \le 2^{c} \cdot 2^{-K(x)}/P(x) \)
  for every integrable test \( t = 2^{d} \) (\texttt{integrableTest\_le\_canonical}). Here
  \( P t \) is a constructive semimeasure and the coding theorem applies.
  \texttt{universalSemimeasure\_div\_isUniversal} states the theorem with \( \m(x)/P(x) \), for
  any universal semimeasure \( \m \).
\item With his constant: \( 2^{-K(y_{d})} 2^{-K(y_{P})} t(x) \le 2^{c} \cdot 2^{-K(x)}/P(x) \)
  for descriptions \( y_{d} \) of \( t \) and \( y_{P} \) of \( P \), where \( c \) depends on
  neither \( P \), \( t \) nor \( x \) (\texttt{integrableTest\_le\_canonical\_exact}; the
  exponent form \( t(x) \le 2^{c + K(y_{d}) + K(y_{P})} \cdot 2^{-K(x)}/P(x) \) is
  \texttt{integrableTest\_le\_canonical\_exact'}). Instead of~\eqref{eq:m(nu)}
  and~\eqref{eq:exact-domin} the proof applies Theorem~\ref{thm:ch2-exact-domination} to the
  product \( G(x, \ang{y,y'}) = f_{y}(x) f_{y'}(x) \) of the described functions.
\end{itemize}
The test is lower semicomputable because \( K \) is upper semicomputable and \( P \) is
computable, so that \( 1/P \) is lower semicomputable. Its expectation is at most
\( \sum_{x} 2^{-K(x)} \le 1 \) by Kraft's inequality.
\end{formalization}

The simple form of our universal test suggests remarkable interpretations.
It says that the outcome $x$ is random with respect to the distribution
$P$, if the latter assigns to $x$ a large enough
probability---however, not in absolute terms, only relatively to the
universal semimeasure $\m(x)$.
The relativization is essential, since
otherwise we could not distinguish between random and nonrandom
outcomes for the uniform distribution $P_{n}$ defined in
Section~\ref{sec:tests}.

For any (not necessarily computable) distribution $P$, the
$P$-expected value of the function $\m(x)/P(x)$ is at most 1,
therefore the relation
 \[
 	\m(x) \le k P(x)
 \]
  holds with probability not smaller than $1-1/k$.
If $P$ is computable then the inequality
 \[
	\m(P) P(x) \lem \m(x)
 \]
  holds for all $x$.
Therefore the following applications are obtained.

\begin{itemize}

 \item If we assume $x$ to be the outcome of an experiment with some
simple computable probability distribution $P$ then $\m(x)$ is a good
estimate of $P(x)$.
The goodness depends on how simple $P$ is to
define and how random $x$ is with respect to $P$: how justified
is the assumption.
Of course, we cannot compute $\m(x)$ but it is
nevertheless well defined.

 \item If we trust that a given object $x$ is random with respect to a
given distribution $P$ then we can use $P(x)$ as an estimate of $\m(x)$.
The degree of approximation depends on the same two factors.

 \end{itemize}

Let us illustrate the behavior of the universal semimeasure $\m(x)$
when $x$ runs over the set of natural numbers.
We define the
semimeasures $v$ and $w$ as follows.
Let $v(n)=c/n^2$ for an
appropriate constant $c$, let $w(n)=v(x)$ if $n=2^{2^k}$, and 0
otherwise.
Then $\m(n)$ dominates both $v(n)$ and $w(n)$.
The function $\m(n)$ dominates $1/n\log^2n$, but jumps at many places
higher than that.
Indeed, it is easy to see that $\m(n)$ converges to
0 slower than any positive computable function converging to 0: in
particular, it is not computable.
Hence it cannot be a measure, that is $\sum_{x} \m(x) < 1$.
We feel that we can make $\m(x)$ ``large'' whenever $x$ is ``simple''.

We cannot compare the universal test defined in the present section
directly with the one defined in Section~\ref{sec:tests}.
There, we investigated tests for a whole family $P_{n}$ of
distributions, where $P_{n}$ is the uniform distribution on the set
$\dB^{n}$.
That the domain of these distributions is a different one for
each $n$ is not essential since we can identify the set $\dB^{n}$ with
the set $\{2^{n},\ldots,2^{n+1}-1\}$.
Let us rewrite the function $d_{0}(x)$ defined in~\eqref{eq:d0} as
 \[
 	d_{0}(x) = n - C(x \mvert n)
 \]
 for $x \in \dB^{n}$, and set it $\infty$ for $x$ outside $\dB^{n}$. Now the
similarity to the expression
 \[
  	\ol\d_{P}(x) = -\log P(x) - \log \m(x)
 \]
  is striking because $n = \log P_{n}(x)$.
Together with the remark
made in the previous paragraph, this observation suggests that the
value of $-\log \m(x)$ is close to the complexity $C(x)$.

\subsection{Conservation of randomness}

The idea of conservation of randomness has been expressed first by Levin,
who has published some papers culminating in~\cite{LevinRandCons84}
developing a general theory.
In this section, we address the issue in its simplest form only.
Suppose that we want to talk about the randomness of integers.
One and the same integer $x$
can be represented in decimal, binary, or in some
other notation: this way, every time a different string will represent the
same number.
Still, the form of the representation should not affect the question of
randomness of $x$ significantly.
How to express this formally?

Let $P$ be a computable measure and let $f : \dS \to \dS$ be a computable
function.
We expect that under certain conditions, the randomness of $x$ should imply
the randomness of $f(x)$---but, for which distribution?
For example, if $x$ is a binary string and $f(x)$ is the decimal notation
for the number expressed by $1x$ then we expect $f(x)$ to be random not
with respect to $P$, but with respect to the measure $f^{*}P$ that is
the image of $P$ under $f$.
This measure is defined by the formula
 \begin{equation}\label{eq:measure-image}
   (f^{*}P)(y) = P(f^{-1}(y)) = \sum_{f(x) = y} P(x).
 \end{equation}

% Gacs source line 2722
 \begin{proposition}
  \label{prop:g2.2.5}
  \lean{Kolmogorov.Randomness.pushforward, Kolmogorov.Randomness.pushMass, Kolmogorov.Randomness.pushMass_isLSC₁, Kolmogorov.Randomness.IsLSC₁.isComputable_of_tsum_eq_one}
  \leanok
  \uses{def:ch2-computable-measure}
If $f$ is a computable function and $P$ is a computable measure
then $f^{*}P$ is a computable measure.
 \end{proposition}
 \begin{proof}
  \proves{prop:g2.2.5}
  \leanok
  \uses{prop:constructive-measure}
It is sufficient to show that $f^{*}P$ is lower semicomputable, since we
have seen that a lower semicomputable measure is computable.
But, the semicomputability can be seen immediately from the definition.
 \end{proof}

\begin{formalization}
\texttt{pushforward f hf P} is \( f^{*}P \) as a computable distribution. ``A lower
semicomputable measure is computable'' is Proposition~\ref{prop:constructive-measure},
re-proved here for the two-sided computability of
Definition~\ref{def:ch2-computable-measure}.
\end{formalization}

Let us see that randomness is indeed preserved for the image.
Similarly to~\eqref{eq:m(nu)}, let us define for an arbitrary computable
function $f$:
 \begin{equation}\label{eq:m(f)}
   \m(f) = \sum\setof{\m(p) : \txt{program $p$ computes $f$}}.
 \end{equation}

% Gacs source line 2739
 \begin{theorem}\label{thm:simple-rand-cons-det}
  \lean{Kolmogorov.Randomness.canonical_pushforward_le, Kolmogorov.Randomness.canonical_pushforward_comp_isIntegrableTest, Kolmogorov.Randomness.det_canonical_image_le, Kolmogorov.Randomness.det_canonical_image_le_exact}
  \leanok
  \uses{prop:g2.2.5, def:g2.2.4, def:ch2-description}
 Let $f$ be a computable function.
Between the randomness test for $P$ and
the test for the image of $P$, the following relation holds, for
all $x$:
  \begin{equation}\label{eq:simple-rand-cons-det}
  \ol\d_{f^{*} P }(f(x)) \lea \ol\d_{P}(x) + K(f) + K(P).
  \end{equation}
 \end{theorem}
 \begin{proof}
  \proves{thm:simple-rand-cons-det}
  \leanok
  \uses{thm:univ-integrable-test, thm:ch2-exact-domination, l:kraft}
Let us denote the function on the left-hand side
of~\eqref{eq:simple-rand-cons-det} by
 \[
   d_{P}(x) = \ol\d_{f^{*}P}(f(x)) = \log\frac{\m(f(x))}{(f^{*}P)(f(x))}.
 \]
It is lower semicomputable, with the help of a program of length
$K(f) + K(P)$, by its very definition.
Let us check that it is an integrable test, so it
satisfies the inequality~\eqref{eq:newtestcond}.
We have
 \[
 \sum_{x} P(x) 2^{d_{P}(x)}
     = \sum_{y} (f^{*}P)(y) 2^{\ol\d_{f^{*}P}(y)} =
       \sum_{y} \m(y) \le 1.
 \]
Hence $d_{P}(x)$ is a test, and hence it is
$\lea \ol\d_{P}(x) + K(f) + K(P)$, by the universality of the test $\ol\d_{P}(x)$.
(Strictly speaking, we must modify the proof of
Theorem~\ref{thm:univ-integrable-test} and see that a program of length $K(f) +
K(P)$ is also sufficient to define the semimeasure $P(x) d(x \mvert P)$.)
 \end{proof}

\begin{formalization}
Multiplicatively, \( 2^{\ol\d_{f^{*}P}(f(x))} \le 2^{c} \cdot 2^{\ol\d_{P}(x)} \), with a
constant that depends on \( f \) and \( P \) (\texttt{canonical\_pushforward\_le}). As in his
proof, \( x \mapsto 2^{\ol\d_{f^{*}P}(f(x))} \) is shown to be an integrable test for \( P \)
(\texttt{canonical\_pushforward\_comp\_isIntegrableTest}); \( 2^{-K(f(x))} \) and
\( 2^{-K(x)} \) are never compared pointwise. With his constant \( K(f) + K(P) \) the theorem
is \texttt{det\_canonical\_image\_le\_exact}. There \( f \) is described through its graph
indicator, and \( f^{*}P \) is obtained as the image under the deterministic transition
\( T(x,y) = 1[y = f(x)] \) of Theorem~\ref{thm:g2.2.3}, which has the same values. The weight
\( \m(f) \) of~\eqref{eq:m(f)} is not formalized.
\end{formalization}

The appearance of the term $K(f)$ on the right-hand side is understandable:
using some very complex function $f$, we can certainly turn any string
into any other one, also a random string into a much less random one.
On the other hand, the term $K(P)$ appears only due to the
imperfection of our concepts.
It will disappear in the more advanced theory presented in later sections,
where we develop \emph{uniform tests}.

% Gacs source line 2779
\begin{example}
For each string $x$ of length $n \ge 1$, let
 \[
  P(x) = \frac{2^{-n}}{n(n+1)}.
 \]
This distribution is uniform on strings of equal length.
Let $f(x)$ be the function that erases the first $k$ bits.
Then for $n \ge 1$ and for any string $x$ of length $n$
we have
 \[
  (f^{*}P)(x) = \frac{2^{-n}}{(n+k)(n+k+1)}.
 \]
(The empty string has the rest of the weight of $f^{*}P$.)
The  new distribution is still uniform on strings of equal length,
though
the total probability of strings of length $n$ has changed somewhat.
We certainly expect randomnes conservation here:
after erasing the first $k$ bits of
a random string of length $n+k$, we should
still get a random string of length $n$.
\end{example}

The computable
function $f$ applied to a string $x$ can be viewed as some kind of
transformation $x \mapsto f(x)$.
As we have seen, it does not make $x$ less random.
Suppose now that we introduce randomization into the transformation
process itself: for example, the machine computing $f(x)$ can also toss
some coins.
We want to say that randomness is also conserved under such more general
transformations, but first of all, how to express such a transformation
mathematically?
We will describe it by a ``matrix'':
a computable probability transition function $T(x,y) \ge 0$ with the
property that
 \[
   \sum_{y} T(x, y) = 1.
 \]
Now, the image of the distribution $P$ under this transformation can be
written as $T^{*}P$, and defined as follows:
 \[
   (T^{*}P)(y) = \sum_{x} P(x) T(x, y).
 \]
How to express now randomness conservation?
It is certainly not true that every possible outcome $y$ is as random
with respect to $T^{*}P$ as $x$ is with respect to $P$.
We can only expect that for each $x$, the conditional probability
(in terms of the transition $T(x, y)$) of those $y$ whose non-randomness is
larger than that of $x$, is small.
This will indeed be expressed by the corollary below.
To get there, we upperbound the $T(x, \cdot)$-expected value
of $2^{\ol\d_{T^{*}P}(y)}$.
Let us go over to exponential notation:

% Gacs source line 2833
 \begin{definition}
  \label{def:g2.2.7}
  \lean{Kolmogorov.Randomness.canonicalIntegrableTest}
  \leanok
  \uses{def:g2.2.4}
Denote
 \begin{align}\label{eq:e.ol-t}
    \ol\t_{P}(x) = 2^{\ol\d_{P}(x)}=\frac{\m(x)}{P(x)}.
  \end{align}
 \end{definition}

\begin{formalization}
This exponential form is the one the formalization uses throughout, with \( 2^{-K(x)} \) in
place of \( \m(x) \) (Definition~\ref{def:g2.2.4}).
\end{formalization}

% Gacs source line 2840
 \begin{theorem}
  \label{thm:g2.2.3}
  \lean{Kolmogorov.Randomness.ComputableTransition, Kolmogorov.Randomness.ComputableTransition.image, Kolmogorov.Randomness.transTest, Kolmogorov.Randomness.transTest_isIntegrableTest, Kolmogorov.Randomness.transTest_le_canonical, Kolmogorov.Randomness.transTest_le_canonical_exact}
  \leanok
  \uses{def:g2.2.7, def:g2.2.1, def:ch2-description}
 We have
 \begin{equation}\label{eq:simple-rand-cons-rand}
 \log \sum_{y} T(x, y) 2^{\ol\d_{T^{*} P}(y)} \lea \ol\d_{P}(x) + K(T) + K(P).
 \end{equation}
 \end{theorem}
 \begin{proof}
  \proves{thm:g2.2.3}
  \leanok
  \uses{thm:univ-integrable-test, thm:ch2-exact-domination, prop:constructive-measure, l:kraft}
Let us denote the function on the left-hand side
of~\eqref{eq:simple-rand-cons-rand} by $t_{P}(x)$.
It is lower semicomputable by its very construction, using a program of
length $\lea K(T) + K(P)$.
Let us check that it satisfies $\sum_{x} P(x) t_{P}(x) \le 1$ which,
in this notation, corresponds to inequality~\eqref{eq:newtestcond}.
We have
 \begin{align*}
   \sum_{x} P(x) t_{P}(x) &= \sum_{x} P(x) \sum_{y} T(x, y) t_{T^{*} P}(y)
\\   &= \sum_{x} P(x) \sum_{y} T(x, y) \frac{\m(y)}{(T^{*}P)(y)}
      = \sum_{y} \m(y) \le 1.
 \end{align*}
It follows that $d_{P}(x) = \log\, t_{P}(x)$ is an integrable
test and hence $d_{P}(x) \lea \ol\d_{P}(x) + K(T) + K(P)$.
(See the remark at the end of the proof of
Theorem~\ref{thm:simple-rand-cons-det}.)
 \end{proof}

\begin{formalization}
Multiplicatively, \( \sum_{y} T(x,y)\, \ol\t_{T^{*}P}(y) \le 2^{c}\, \ol\t_{P}(x) \)
(\texttt{transTest\_le\_canonical}), with \( 2^{-K} \) in place of \( \m \). The transition
need only be lower semicomputable in \( (x,y) \), with rows summing to \( 1 \)
(\texttt{ComputableTransition}). Together with the normalization this implies computability,
by the argument of Proposition~\ref{prop:constructive-measure} run row by row (not
formalized), so nothing is lost. The image \( T^{*}P \) is a computable distribution by
Proposition~\ref{prop:constructive-measure}.

With his constant \( K(T) + K(P) \) the theorem is \texttt{transTest\_le\_canonical\_exact}:
\( 2^{-K(y_{T})} 2^{-K(y_{P})} \sum_{y} T(x,y)\, \ol\t_{T^{*}P}(y) \le 2^{c}\, \ol\t_{P}(x) \),
where \( y_{T} \) describes \( T \) on encoded pairs and \( c \) depends on neither \( T \),
\( P \) nor \( x \). His closing remark about a program of length \( K(T) + K(P) \) becomes one
family, lower semicomputable in \( x \) and the index \( \ang{y_{T}, y_{P}} \), to which
Theorem~\ref{thm:ch2-exact-domination} applies. In it, \( 1/(T^{*}P)(y) \) is only upper
semicomputable; since \( T^{*}P \) has total mass \( 1 \) it is written as
\( (1 - \sum_{y' \ne y} (T^{*}P)(y'))^{-1} \), which is lower semicomputable.
\end{formalization}

% Gacs source line 2864
 \begin{corollary}
  \label{cor:g2.2.8}
  \lean{Kolmogorov.Randomness.transition_deficiency_tail, Kolmogorov.Randomness.transition_deficiency_tail_exact}
  \leanok
  \uses{thm:g2.2.3}
There is a constant $c_{0}$
such that for every integer $k \ge 0$, for all $x$ we have
 \[
 \sum\setof{ T(x, y) : \ol\d_{T^{*}P}(y) - \ol\d_{P}(x) > k + K(T) + K(P)}
 \le 2^{-k + c_{0}}.
 \]
 \end{corollary}
 \begin{proof}
  \proves{cor:g2.2.8}
  \leanok
  \uses{thm:g2.2.3}
The theorem says
 \[
  \sum_{y} T(x, y) \ol\t_{T^{*} P}(y) \lem t_{P}(x)2^{K(T) + K(P)}.
 \]
Thus, it upperbounds the expected value of
the function  $2^{\ol\d_{T^{*}P)}(y)}$ according to the distribution $T(x, \cdot)$.
Applying Markov's inequality~\eqref{eq:Markov}
to this function yields the desired result.
 \end{proof}

\begin{formalization}
Multiplicatively, the event is
\( 2^{k}\, \ol\t_{P}(x) < 2^{-K(y_{T})} 2^{-K(y_{P})}\, \ol\t_{T^{*}P}(y) \), and its
\( T(x,\cdot) \)-probability is at most \( 2^{c-k} \) with \( c \) independent of \( T \),
\( P \), \( k \) and \( x \) (\texttt{transition\_deficiency\_tail\_exact}). Without the term
\( K(T) + K(P) \), and with \( c \) depending on \( T \) and \( P \), it is
\texttt{transition\_deficiency\_tail}.
\end{formalization}

\section{Infinite sequences}

These lecture notes treat the theory of randomness over continuous spaces mostly
separately, starting in Section~\ref{sec:cont-spaces}.
But the case of computable measures over infinite sequences is particularly
simple and appealing, so we give some of its results here, even if most
follow from the more general results given later.

In this section, we will fix a finite or countable \df{alphabet}
\( \Sg=\set{s_{1},s_{2},\dots} \), and consider
probability distributions over the set
\begin{align*}
   X = \Sg^{\dN}
 \end{align*}
of infinite sequences with members in \( \Sg \).
An alternative way of speaking of this is to consider a sequence of
\df{random variables} \( X_{1},X_{2},\dots \), where \( X_{i}\in\Sg \), with a joint
distribution.
The two interesting extreme special cases are \( \Sg=\dB=\{0,1\} \), giving the set
of infinite 0-1 sequences, and \( \Sg=\dN \), giving the set sequences of natural
numbers.

\begin{formalization}
The formalization takes \( \Sg = \dB \): a sequence is a function \( \dN \to \{0,1\} \)
(\texttt{Seq}), a string is a list of bits. Everything is stated through strings: open sets
by generating sets of strings, measures by their values on cylinders, lower semicomputable
functions on \( X \) by their values on prefixes. His general measure theory,
Definitions~\ref{def:g2.3.15}--\ref{def:g2.3.27} with their propositions, is not formalized
(see the note at the start of Subsection 2.3.2).
\end{formalization}

\subsection{Null sets}

Our goal is here to illustrate the measure theory developed in
Section~A.2, through Subsection~A.2.3,
on the concrete example of the set
of infinite sequences, and then to develop the theory of randomness
in it.

We will distinguish a few simple kinds of subsets of the set \( \S \) of sequences,
those for which probability can be defined especially easily.

% Gacs source line 2917
\begin{definition}\
  \label{def:g2.3.1}
  \lean{Kolmogorov.Randomness.Seq, Kolmogorov.Randomness.cylinder, Kolmogorov.Randomness.openOf}
  \leanok
  \begin{bullets}
  \item
  For string \( z\in\Sg^{*} \), we will denote by \( zX \) the set of elements of \( X \)
with prefix \( z \).
Such sets will be called \df{cylinder sets}.
 \item
A subset of \( X \) is \df{open} if it is the union of any number (not necessarily
finite) of cylinder sets.
It is called \df{closed} if its complement is open.
 \item
An open set is called \df{constructive} if it is the union of a recursively
enumerable set of cylinder sets.
\item A set is called \( G_{\dg} \) if it is the intersection of a sequence of open sets.
 It is called \( F_{\sg} \) if it is the union of a sequence of closed sets.
\item
 We say that a set \( E\sbsq X \) is \df{finitely determined} if there is an \( n \)
and an \( \cE\sbsq \Sg^{n} \) such that
\begin{align*}
  E = \cE X=\bigcup_{s\in\cE}s X.
 \end{align*}
Let \( \cF \) be the class of all finitely determined subsets of \( X \).
  \item
A class of subsets of \( X \) is called an \df{algebra} if it is closed with respect
to finite intersections and complements (and then of course, also with respect
to finite unions).
It is called a \( \sg \)-algebra (sigma-algebra)
when it is also closed with respect to countable
intersections (and then of course, also with respect to countable unions).
\end{bullets}
\end{definition}

\begin{formalization}
Formalized for \( \Sg = \dB \): the cylinder \( zX \) is \texttt{cylinder z}, and the open set
generated by a set \( \Gg \) of strings is \texttt{openOf \( \Gg \)}. A constructive open set is
\texttt{openOf \( \Gg \)} with \( \Gg \) recursively enumerable, which is how
Definition~\ref{def:g2.3.13} uses it. Mathlib's product topology on \( \dN \to \{0,1\} \) has
these open sets; the formalization uses only that cylinders are open, closed and compact.
\( G_{\dg} \) and \( F_{\sg} \) sets, finitely determined sets and algebras are not restated.
\end{formalization}

% Gacs source line 2949
\begin{example}
An example open set that is not finitely determined, is the set \( E \) of all sequences that
contain a substring \( 11 \).
\end{example}

The following observations are easy to prove.

% Gacs source line 2956
\begin{proposition}\
  \label{prop:g2.3.3}
  \lean{Kolmogorov.Randomness.isOpen_cylinder, Kolmogorov.Randomness.isClosed_cylinder, Kolmogorov.Randomness.isOpen_openOf}
  \uses{def:g2.3.1}
  \begin{alphenum}
 \item
 Every open set \( G \) can be represented as the union of a sequence of \emph{disjoint}
cylinder sets.
 \item
The set of open sets is closed with respect to finite intersection
and arbitrarily large union.
  \item
 Each finitely determined set is both open and closed.
 \item
The class \( \cF \) of finitely determined sets forms an algebra.
 \end{alphenum}
\end{proposition}

\begin{formalization}
Partially formalized: cylinders are open and closed, and \texttt{openOf \( \Gg \)} is open.
Clause~(b) holds in any topology. Clause~(a) is not needed, because the measure of an open
set is defined without a disjoint representation (Definition~\ref{def:g2.3.8}).
Clauses~(c) and~(d) are not formalized.
\end{formalization}

Probability is generally defined as a nonnegative, monotonic
function (a ``measure'', see later) on some subsets of the event
space (in our case the set \( X \)) that has a certain additivity property.
We do not need immediately the complete definition of measures, but let us say
what we mean by an additive set function.

% Gacs source line 2977
\begin{definition}
  \label{def:g2.3.4}
  \notready
Consider a nonnegative, function \( \mu \) is defined over some subsets of \( X \).
that is also \df{monotonic}, that is \( A\sbsq B \) implies  \( \mu(A)\sbsq\mu(B) \).
We say that it is \df{additive} if,
whenever it is defined on the sets
\( E_{1},\dots,E_{n} \), and on \( E=\bigcup_{i=1}^{n}E_{i} \)
and the sets \( E_{i} \) are mutually disjoint, then we have
\( \mu(E)=\mu(E_{1})+\dots+\mu(E_{n}) \).

We say that \( \mu \) is \df{countably additive}, or \( \sg \)-additive (sigma-additive),
if in addition whenever it is defined on the sets
\( E_{1},E_{2},\dots \), and on \( E=\bigcup_{i=1}^{\infty}E_{i} \)
and the sets \( E_{i} \) are mutually disjoint, then we have
\( \mu(E)=\sum_{i=1}^{\infty}\mu(E_{i}) \).
\end{definition}

\begin{formalization}
Not formalized. This is the general notion of an additive set function, which Mathlib has
(\texttt{MeasureTheory.AddContent}, \texttt{MeasureTheory.Measure}). The formalization uses
additivity only on cylinders, in the form of Definition~\ref{def:g2.3.5}.
\end{formalization}

First we consider only measures defined on cylinder sets.

% Gacs source line 2995
\begin{definition}[Measure over \( \Sg^{\dN} \)]
  \label{def:g2.3.5}
  \lean{Kolmogorov.Randomness.IsSeqMeasure, Kolmogorov.Randomness.ComputableSeqMeasure, Kolmogorov.Randomness.uniformSeqMeasure}
  \leanok
  \uses{def:g2.3.1, def:ch2-computable-measure}
Let \( \mu:\Sg^{*}\to\dR_{+} \) be a function assigning a nonnegative number to
each string in \( \Sg^{*} \).
We will call \( \mu \) a \df{measure} if it satisfies the condition
 \begin{align}\label{eq:extension-sum}
   \mu(x)=\sum_{s\in\Sg}\mu(xs).
 \end{align}
We will take the liberty to also write
 \begin{align*}
   \mu(s X)=\mu(s)
 \end{align*}
for all \( s\in\Sg^{*} \), and \( \mu(\emptyset)=0 \).
We call \( \mu \) a \df{probability measure} if \( \mu(\Lg)=1 \) and correspondingly,
\( \mu(X)=1 \).

A measure is called \df{computable} if it is computable as a function
\( \mu:\Sg^{*}\to\dR_{+} \) according to Definition~\ref{def:semicomputability}.
\end{definition}

\begin{formalization}
Formalized for \( \Sg = \dB \): \texttt{IsSeqMeasure \( \mu \)} is
\( \mu(x) = \mu(x0) + \mu(x1) \) for a function from strings to \( [0,\infty] \).
\texttt{ComputableSeqMeasure} bundles this law with \( \mu(\Lg) = 1 \) and the two-sided
computability of Definition~\ref{def:ch2-computable-measure}; every result about sequences
assumes it. The uniform measure \( \lg(x) = 2^{-\len{x}} \) of
Examples~\ref{example:null-set} is \texttt{uniformSeqMeasure}.
\end{formalization}

The following observation lets us extend measures.

% Gacs source line 3016
\begin{lemma}\label{lem:cyl-measure}
  \lean{Kolmogorov.Randomness.openMeasure_mono_openOf, Kolmogorov.Randomness.exists_level_of_cylinder_subset}
  \leanok
  \uses{def:g2.3.5, def:g2.3.8}
  Let \( \mu \) be a measure defined as above.
If a cylinder set \( s X \) is the disjoint union \( xX=x_{1}X\cup x_{2}X\cup\dots \)
of cylinder sets, then we have \( \mu(z)=\sum_{i}\mu(x_{i}) \).
\end{lemma}
\begin{proof}
  \proves{lem:cyl-measure}
  \leanok
  \uses{def:g2.3.1}
Let us call an arbitrary \( y\in\Sg^{*} \) \df{good} if
 \begin{align*}
   \mu(zy)=\sum_{i}\mu(zyX\cap x_{i}X)
 \end{align*}
holds.
We have to show that the empty string \( \Lg \) is good.

It is easy to see that \( y \) is good if \( zyX\sbsq x_{i}X \) for some \( i \).
Also, if \( ys \) is good for all \( s\in\Sg \) then \( y \) is good.
Now assume that \( \Lg \) is not good.
Then there is also an \( s_{1}\in\Sg \) that is not good.
But then there is also an \( s_{2}\in\Sg \) such that \( s_{1}s_{2} \) is not good.
And so on, we obtain an infinite sequence \( s_{1}s_{2}\dotsm \) such that
\( s_{1}\dotsm s_{n} \) is not good for any \( n \).
But then the sequence \( z s_{1}s_{2}\dotsm \) is not contained in
\( \bigcup_{i}x_{i}X \), contrary to the assumption .
\end{proof}

\begin{formalization}
Formalized as monotonicity of the measure in the open set: if the open set generated by
\( \Gg \) is contained in the one generated by \( \Gg' \), then \( \mu(\Gg) \le \mu(\Gg') \)
for the measure of Definition~\ref{def:g2.3.8}. So the measure depends only on the open set,
not on its generators, which is what the lemma is for. The compactness step of his proof
(the infinite sequence \( s_{1}s_{2}\dotsm \)) is taken from Mathlib: \( \dN \to \{0,1\} \) is
compact, so a cylinder covered by cylinders is covered by finitely many of them, at a common
length (\texttt{exists\_level\_of\_cylinder\_subset}).
\end{formalization}

% Gacs source line 3040
\begin{corollary}
  \label{cor:g2.3.7}
  \lean{Kolmogorov.Randomness.openMeasure_le_tsum, Kolmogorov.Randomness.sum_le_levelMass_of_prefixFree}
  \uses{lem:cyl-measure, def:g2.3.8}
  Two disjoint union representations of the same open set
 \begin{align*}
   \bigcup_{i}x_{i}X = \bigcup_{i}y_{i}X
 \end{align*}
imply \( \sum_{i}\mu(x_{i})=\sum_{i}\mu(y_{i}) \).
\end{corollary}
\begin{proof}
  \proves{cor:g2.3.7}
  \uses{lem:cyl-measure}
  The set \( \bigcup_{i}x_{i}X \) can also be written as
 \begin{align*}
  \bigcup_{i}x_{i}X=\bigcup_{i,j} x_{i}X\cap y_{j}X.
\end{align*}
An element \( x_{i}X\cap y_{j}X \) is nonempty only if one of the strings
\( x_{i},y_{j} \) is a continuation of the other.
Calling this string \( z_{ij} \) we have \( x_{i}X\cap y_{j}X=z_{ij}X \).
Now Lemma~\ref{lem:cyl-measure} is applicable to the union
\( x_{i}X=\bigcup_{j}(x_{i}X\cap y_{j}X) \),
giving
 \begin{align*}
        \mu(x_{i}) &=\sum_{j}\mu(x_{i}X\cap y_{j}X),
\\   \sum_{i}\mu(x_{i}) &=\sum_{i,j}\mu(x_{i}X\cap y_{j}X).
\end{align*}
The right-hand side is also equal similarly to \( \sum_{j}\mu(y_{j}) \).
\end{proof}

\begin{formalization}
Not stated as an equality in Lean; this is listed as open in the strategy file. Both
inequalities are there. The measure of the open set generated by \( \Gg \) is at most
\( \sum_{x \in \Gg} \mu(x) \) (\texttt{openMeasure\_le\_tsum}), and a finite prefix-free set
has total measure at most the level measure of the open set it generates
(\texttt{sum\_le\_levelMass\_of\_prefixFree}). For prefix-free \( \Gg \) they give
\( \mu(\Gg) = \sum_{x \in \Gg} \mu(x) \), and Lemma~\ref{lem:cyl-measure} then gives the
corollary.
\end{formalization}

The above corollary allows the following definition:

% Gacs source line 3067
\begin{definition}
  \label{def:g2.3.8}
  \lean{Kolmogorov.Randomness.openMeasure, Kolmogorov.Randomness.levelMass}
  \leanok
  \uses{def:g2.3.5, def:g2.3.1}
For an open set \( G \) given as a disjoint union of cylinder sets \( x_{i}X \)
let \( \mu(G)=\sum_{i}\mu(x_{i}) \).
For a closed set \( F=X\xcpt G \) where \( G \) is open, let
\( \mu(F)=\mu(X)-\mu(G) \).
\end{definition}

\begin{formalization}
Defined directly, without a disjoint representation: \texttt{openMeasure \( \mu \) \( \Gg \)} is
the supremum over \( n \) of the total measure of the strings of length \( n \) that have a
prefix in \( \Gg \) (\texttt{levelMass}). For a disjoint representation this is
\( \sum_{i} \mu(x_{i}) \) (Corollary~\ref{cor:g2.3.7}). The measure of a closed set is not
formalized.
\end{formalization}

The following is easy to check.

% Gacs source line 3076
\begin{proposition}
  \label{prop:g2.3.9}
  \lean{Kolmogorov.Randomness.openMeasure_mono_openOf, Kolmogorov.Randomness.openMeasure_mono_set, Kolmogorov.Randomness.openMeasure_iUnion_le}
  \uses{def:g2.3.8}
The measure as defined on open sets is monotonic and
countably additive.
In particular, it is countably additive
on the algebra of finitely determined sets.
\end{proposition}

\begin{formalization}
Partially formalized: monotonicity (\texttt{openMeasure\_mono\_openOf}) and countable
subadditivity \( \mu(\bigcup_{i} G_{i}) \le \sum_{i} \mu(G_{i}) \)
(\texttt{openMeasure\_iUnion\_le}), which is all the chapter uses. Additivity over disjoint
open sets and the clause on finitely determined sets are not formalized.
\end{formalization}

Before extending measures to a wider range of sets, consider an important
special case.
With infinite sequences, there are certain events that are not impossible, but
still have probability 0.

% Gacs source line 3088
\begin{definition}
  \label{def:g2.3.10}
  \lean{Kolmogorov.Randomness.IsNullSet}
  \leanok
  \uses{def:g2.3.8}
Given a measure \( \mu \), a set \( N\sbsq X \) is called a \df{null set} with respect
to \( \mu \)
if there is a sequence \( G_{1},G_{2},\dots \) of open sets with the property
\( N\sbsq \bigcap_{m}G_{m} \), and \( \mu(G_{m})\le 2^{-m} \).
We will say that the sequence \( G_{m} \) \df{witnesses} the fact that \( N \) is a null set.
\end{definition}

% Gacs source line 3096
\begin{examples}\label{example:null-set}
Let \( \Sg=\{0,1\} \), let us define the measure \( \lg \) by
\( \lg(x)=2^{-n} \) for all \( n \) and for all \( x\in\Sg^{n} \).
  \begin{enumerate}
  \item\label{i:singleton}
For every infinite sequence \( \xi\in X \), the one-element set \( \{\xi\} \) is a null set
with respect to \( \lg \).
Indeed, for each natural number \( n \), let
\( H_{n}=\setof{\xi\in X: \xi(0)=\xi(0),\xi(1)=\xi(1),\dots,\xi(n)=\xi(n)} \).
Then \( \lg(H_{n})=2^{-n-1} \), and \( \{s\}=\bigcap_{n}H_{n} \).
\item\label{i:even-places}
For a less trivial example,
consider the set \( E \) of those elements \( t\in X \) that are 1 in each positive even
position, that is
 \begin{align*}
   E=\setof{t\in X: t(2)=1,t(4)=1,t(6)=1\dots}.
 \end{align*}
Then \( E \) is a null set.
Indeed, for each natural number \( n \), let
\( G_{n}=\setof{t\in X: t(2)=1,t(4)=1,\dots,t(2n)=1} \).
This helps expressing \( E \) as \( E=\bigcap_{n}G_{n} \), where \( \lg(G_{n})=2^{-n} \).
  \end{enumerate}
\end{examples}

% Gacs source line 3120
\begin{proposition}\label{propo:infin-null-sets}
  \lean{Kolmogorov.Randomness.isNullSet_iUnion}
  \leanok
  \uses{def:g2.3.10}
Let \( N_{1},N_{2},\dots \) be a sequence of null sets with respect to a measure
\( \mu \).
Their union \( N=\bigcup_{i}N_{i} \) is also a null set.
\end{proposition}
\begin{proof}
  \proves{propo:infin-null-sets}
  \leanok
  \uses{prop:g2.3.9}
  Let \( G_{i,1},G_{i,2},\dots \) be the infinite sequence witnessing the fact that
\( N_{i} \) is a null set.
Let \( H_{m}=\bigcup_{i=1}^{\infty} G_{i,m+i} \).
Then the sequence \( H_{m} \) of open sets witnesses the fact that \( N \) is a null set.
\end{proof}

We would like to extend our measure to null sets \( N \) and say \( \mu(N)=0 \).
The proposition we have just proved shows that such an extension would be
countably additive on the null sets.
But we do not need this extension for the moment, so we postpone it.
Still, given a probability measure \( P \) over \( X \), it becomes meaningful to say
that a certain property holds with probability 1.
What this means is that the set of those sequences that do not have this
property is a null set.

Following the idea of Martin-L\"of, we would like to call a sequence
nonrandom, if it is contained in some \emph{simple} null set; that is it has
some easily definable property with probability 0.
As we have seen  in part~\ref{i:singleton} of Example~\ref{example:null-set},
it is important to insist on simplicity, otherwise (say
with respect to the measure \( \lg \)) every sequence might be contained in a null set,
namely the one-element set consisting of itself.
But most of these sets are not defined simply at all.
An example of a simply defined null set is given is
part~\ref{i:even-places} of Example~\ref{example:null-set}.
These reflections justify the following definition, in which ``simple'' is
specified as ``constructive''.

% Gacs source line 3154
\begin{definition}
  \label{def:g2.3.13}
  \lean{Kolmogorov.Randomness.IsConstructiveNull}
  \leanok
  \uses{def:g2.3.10, def:g2.3.5}
Let \( \mu \) be a computable measure.
A set \( N\sbsq X \) is called a \df{constructive null set} if there is
recursively enumerable set \( \Gg\sbsq\dN\times\Sg^{*} \)
with the property that denoting \( \Gg_{m}=\setof{x: \tup{m,x}\in\Gg} \)
and \( G_{m}=\bigcup_{x\in\Gg_{m}}xX \)
we have \( N\sbsq\bigcap_{m}G_{m} \), and \( \mu(G_{m})\le 2^{-m} \).
\end{definition}

\begin{formalization}
Recursive enumerability of \( \Gg \) is that of the relation \( x \in \Gg_{m} \) in the pair
\( (m,x) \). The definition is stated for any \( \mu \); the theorems assume \( \mu \)
computable.
\end{formalization}

In words, the difference between the definition of null sets and constructive
null sets is that the sets \( G_{m}=\bigcup_{x\in\Gg_{m}}xX \) here
are required to be constructive open, moreover, in such a way that from \( m \) one
can compute the program generating \( G_{m} \).
In even looser words, a set is a constructive null set if for any \( \eps>0 \) one
can construct effectively
a union of cylinder sets containing it, with total measure \( \le\eps \).

Now we are in a position to define random infinite sequences.

% Gacs source line 3173
\begin{definition}[Random sequence]
  \label{def:g2.3.14}
  \lean{Kolmogorov.Randomness.IsRandomSeq}
  \leanok
  \uses{def:g2.3.13}
An infinite sequence \( \xi \) is \df{random} with respect to a probability measure
\( P \) if and only if \( \xi \) is not contained in any constructive null set with
respect to \( P \).
\end{definition}

It is easy to see that the set of nonrandom sequences is a null set.
Indeed, there is only a countable number or constructive null sets,
so even their union is a null set.
The following theorem strengthens this observation significantly.

% Gacs source line 3184
\begin{theorem}\label{thm:univ-seq-test}
  \lean{Kolmogorov.Randomness.isConstructiveNull_universalNull, Kolmogorov.Randomness.subset_universalNull, Kolmogorov.Randomness.isRandomSeq_iff}
  \leanok
  \uses{def:g2.3.14, def:g2.3.5}
Let us fix a \emph{computable} probability measure \( P \).
The set of all nonrandom sequences is a constructive null set.
\end{theorem}
Thus, there is a \df{universal} constructive null set, containing all other
constructive null sets.
A sequence is random when it does not belong to this set.
\begin{proof}[Proof of Theorem~\protect\ref{thm:univ-seq-test}]
  \proves{thm:univ-seq-test}
  \leanok
  \uses{thm:g2.3.2, thm:univ-integrable-test-seq, thm:defic-charac-cpt-seqs}
The proof uses the projection technique that has appeared several times in this
book, for example in proving the existence of a universal lower semicomputable
semimeasure.

We know that it is possible to list all recursively enumerable subsets of the
set \( \dN\times\Sg^{*} \) into a sequence, namely that there is a recursively
enumerable set \( \Dg\sbsq\dN^{2}\times\Sg^{*} \)
with the property that for every recursively enumerable
set \( \Gg\sbsq\dN\times\Sg^{*} \) there is an \( e \) with
\( \Gg=\setof{\tup{m,x}: \tup{e,m,x}\in\Dg} \).
We will write
 \begin{align*}
  \Dg_{e,m}&=\setof{x:\tup{e,m,x}\in\Dg},
\\ D_{e,m}&=\bigcup_{x\in\Dg_{e,m}}xX.
\end{align*}
 We transform the set \( \Dg \) into another set \( \Dg' \)
with the following property.
\begin{bullets}
\item For each \( e,m \) we have \( \sum_{\tup{e,m,x}\in\Dg'}\mu(x) \le 2^{-m+1} \).
\item If  for some \( e \) we have
\( \sum_{\tup{e,m,x}\in\Dg}\mu(x)\le 2^{-m} \) for all \( m \) then for all \( m \) we have
\( \setof{x: \tup{e,m,x}\in\Dg'}=\setof{x: \tup{e,m,x}\in\Dg} \).
\end{bullets}
This transformation is routine, so we leave it to the reader.
By the construction of \( \Dg' \), for every constructive null set \( N \) there is an \( e \)
with
 \begin{align*}
   N\sbsq\bigcap_{m}D'_{e,m}.
 \end{align*}
Define the recursively enumerable set
 \begin{align*}
   \hat\Gg=\setof{\tup{m,x}: \exist{e} \tup{e,m+e+2,x}\in\Dg'}.
 \end{align*}
Then \( \hat G_{m}=\bigcup_{e}D'_{e,m+e+2} \).
For all \( m \) we have
 \begin{align*}
\sum_{x\in\hat\Gg_{m}}\mu(x)&=\sum_{e}\sum_{\tup{e,m+e+2,x}\in\Dg'}\mu(x)
 \le \sum_{e}2^{-m-e-1} = 2^{-m}.
 \end{align*}
This shows that \( \hat\Gg \) defines a constructive null set.
Let \( \Gg \) be any other recursively enumerable subset of \( \dN\times\Sg^{*} \) that
defines a constructive null set.
Then there is an \( e \) such that for all \( m \) we have \( G_{m}=D'_{e,m} \).
The universality follows now from
 \begin{align*}
\bigcap_{m}G_{m}=\bigcap_{m}D'_{e,m}\sbsq \bigcap_{m}D'_{e,m+e+2} \sbsq\bigcap_{m}\hat G_{m}.
 \end{align*}
\end{proof}

\begin{formalization}
The Lean proof does not use the projection technique, so it needs no transformation of
\( \Dg \) into \( \Dg' \). It names the universal constructive null set:
\( \setof{\xi : R(\xi) = \infty} \), where
\( R(\xi) = \sup_{n} 2^{-K(\xi_{1:n})}/P(\xi_{1:n}) \) is the multiplicative form of the
right-hand side of~\eqref{eq:defic-charac-seq}. This set is a constructive null set by
Theorem~\ref{thm:g2.3.2}(a), applied to the universal test of
Theorem~\ref{thm:univ-integrable-test-seq}. A constructive null set \( N \) lies in
\( \setof{\xi : d(\xi) = \infty} \) for an integrable test \( d \)
(Theorem~\ref{thm:g2.3.2}(b)), and \( 2^{d} \le 2^{c} R \) by the second half of the proof of
Theorem~\ref{thm:defic-charac-cpt-seqs}, so \( N \sbsq \setof{R = \infty} \). Hence \( \xi \)
is random if and only if \( R(\xi) < \infty \) (\texttt{isRandomSeq\_iff}).
\end{formalization}

\subsection{Probability space}

\begin{formalization}
Not formalized: Definitions~\ref{def:g2.3.15}--\ref{def:g2.3.27} and
Propositions~\ref{prop:g2.3.19}, \ref{prop:g2.3.21}, \ref{prop:g2.3.24},
\ref{propo:Riesz-extension-seqs} and~\ref{prop:g2.3.28}. This is textbook measure and
integration theory, which Mathlib provides in general form (\texttt{MeasurableSpace},
\texttt{borel}, \texttt{MeasureTheory.Measure}, Lebesgue measure on \( \dR \),
\texttt{Measurable}, \texttt{MeasureTheory.lintegral}), and no result of the chapter needs
it: every statement about sequences is proved through cylinders. In place of integrals,
Definition~\ref{def:g2.3.29} uses level sums over strings of a fixed length. A bridge from
the string-level measures of Definition~\ref{def:g2.3.5} to Mathlib measures on
\( \dN \to \{0,1\} \) is an optional extension.
\end{formalization}

Now we are ready to extend measure to a much wider range of subsets of \( X \).

% Gacs source line 3245
\begin{definition}
  \label{def:g2.3.15}
  \notready
 Elements of the smallest \( \sg \)-algebra \( \cA \) containing the cylinder sets of \( X \) are
called the \df{Borel} sets of \( X \).
The pair \( (X,\cA) \) is an example of a \df{measureable space}, where the Borel
sets are called the \df{measureable sets}.

A nonnegative sigma-additive function \( \mu \) over \( \cA \) is called a \df{measure}.
It is called a \df{probability measure} if \( \mu(X)=1 \).
If \( \mu \) is fixed then the triple \( (X,\cA,\mu) \) is called a \df{measure space}.
If it is a probability measure then the space is called a \df{probability space}.
\end{definition}

In the Appendix we cited a central theorem of measure theory,
Caratheodory's extension theorem.
It implies that if a measure is defined on an algebra \( \cL \) in a sigma-additive way
then it can be extended uniquely
to the \( \sg \)-algebra generated by \( \cL \), that is the
smallest \( \sg \)-algebra containing \( \cL \).
We defined a measure \( \mu \) over \( X \) as a nonnegative function
\( \mu:\Sg^{*}\to\dR_{+} \)
satisfying the equality~\eqref{eq:extension-sum}.
Then we defined \( \mu(xX)=\mu(x) \), and further
extended \( \mu \) in a \( \sg \)-additive way to all elements of the algebra \( \cF \).
Now Caratheodory's theorem allows us to extend it uniquely to all Borel sets,
and thus to define a measureable space \( (X,\cA,\mu) \).

Of course, all null sets in \( \cA \) get measure 0.

Open,  closed and measureable sets can also be defined in the set of real numbers.

% Gacs source line 3275
\begin{definition}
  \label{def:g2.3.16}
  \notready
  A subset \( G\sbsq\dR \) is \df{open} if it is the union of a set of open
  intervals \( \opint{a_{i}}{b_{i}} \).
It is \df{closed} if its complement is open.

The set \( \cB \) of \df{Borel sets} of \( \dB \) is defined as the smallest
\( \sg \)-algebra containing all open sets.
\end{definition}

The pair \( (\dR,\cB) \) is another example of a measureable space.

% Gacs source line 3286
\begin{definition}[Lebesgue measure]
  \label{def:g2.3.17}
  \notready
Consider the set left-closed intervals of the line (including intervals
of the form \( \opint{-\infty}{a} \).
Let \( \cL \) be the set of finite disjoint unions of such intervals.
This is an algebra.
We define the function \( \lg \) over \( \cL \) as follows:
\( \lg(\bigcup_{i} \lint{a_{i}}{b_{i}}) = \sum_{i} b_{i} - a_{i} \).
It is easy to check that this is a \( \sg \)-additive measure and therefore by
Caratheodory's theorem can be extended to the set \( \cB \) of all Borel sets.
This function is called the \df{Lebesgue measure} over \( \dR \), giving us the
measureable space \( (\dR,\cB,\lg) \).
\end{definition}

Finally, we can define the notion of a measureable function over \( X \).

% Gacs source line 3301
\begin{definition}[Measureable functions]
  \label{def:g2.3.18}
  \notready
A function \( f:X\to \dR \) is called \df{measureable} if and only if
\( f^{-1}(E)\in\cA \) for all \( E\in\cB \).
\end{definition}

The following is easy to prove.

% Gacs source line 3308
\begin{proposition}
  \label{prop:g2.3.19}
  \notready
Function \( f:X\to\dR \) is measureable if and only if all sets of the form
\( f^{-1}(\opint{r}{\infty})=\setof{x: f(x)>r} \) are measureable, where \( r \) is a
rational number.
\end{proposition}

\subsection{Computability}

\begin{formalization}
Not formalized (see the note at the start of Subsection 2.3.2). Lower semicomputable
functions on \( X \) enter the formalization only as tests, through their values on prefixes
(Definition~\ref{def:g2.3.29}).
\end{formalization}

Measureable functions are quite general; it is worth introducing some more
resticted kinds of function over the set of sequences.

% Gacs source line 3320
\begin{definition}
  \label{def:g2.3.20}
  \notready
  A function \( f:X\to\dR \) is \df{continuous} if and only if for every \( x\in X \), for
every \( \eps>0 \) there is a cylinder set \( C\ni x \) such that \( |f(y)-f(x)|<\eps \) for all
\( y\in C \).

  A function \( f:X\to\dR \) is \df{lower semicontinuous} if for every \( r\in\dR \) the set
\( \setof{x\in X: f(x)>r} \) is open.
 It is \df{upper semicontinuous} if \( -f \) is lower semicontinuous.
\end{definition}

The following is easy to verify.

% Gacs source line 3332
\begin{proposition}
  \label{prop:g2.3.21}
  \notready
A function \( f:X\to\dR \) is continuous if and only if it is both upper and lower semicontinuous.
\end{proposition}

% Gacs source line 3336
\begin{definition}
  \label{def:g2.3.22}
  \notready
  A function \( f:X\to\dR \) is \df{computable} if for every open rational interval
\( \opint{u}{v} \) the set \( f^{-1}(\opint{u}{v}) \) is a constructive open set of
\( X \), uniformly in \( u,v \).
\end{definition}

Informally this means that if for the infinite sequence \( \xi\in X \) we have
\( u<f(\xi)<v \) then from \( u,v \) sooner or later we will find a prefix \( x \)
of \( \xi \) with the property \( u<f(xX)<v \).

% Gacs source line 3346
\begin{definition}
  \label{def:g2.3.23}
  \notready
A function \( f:X\to\dR \) is \df{lower semicomputable}
if for every rational \( r \) the set \( \setof{s\in X: f(s)>r} \)
is a constructive open set of \( X \), uniformly in \( r \).
It is \df{upper semicomputable} if \( -f \) is lower semicomputable.
\end{definition}

The following is easy to verify.

% Gacs source line 3355
\begin{proposition}
  \label{prop:g2.3.24}
  \notready
  A function \( X\to\dR \) is computable if and only if it is both lower and upper
  semicomputable.
\end{proposition}

\subsection{Integral}

\begin{formalization}
Not formalized (see the note at the start of Subsection 2.3.2).
\end{formalization}

The definition of integral over a measure space is given in the
Appendix, in Subsection~A.2.3.
Here, we give an exposition specialized to infinite sequences.

% Gacs source line 3366
 \begin{definition}
  \label{def:g2.3.25}
  \notready
A measurable function \( f : X \to \dR \) is called a \df{step function} if
its range is finite.
The set of step functions will be called \( \cE \).

Given a step function \( f \) which takes values \( x_{i} \) on sets \( A_{i} \), and a
finite measure \( \mu \), we define
 \[
 \mu(f) = \mu f = \int f\,d\mu = \int f(x) \mu(d x)
   = \sum_{i} x_{i} \mu(A_{i}).
 \]
 \end{definition}

Proposition~A.2.14, when specified to our situation here, says
the following.

% Gacs source line 3382
 \begin{proposition}\label{propo:Riesz-extension-seqs}
  \notready
The functional \( \mu \) defined above on step functions
can be extended to the set
\( \cE_{+} \) of monotonic limits of nonnegative elements of \( \cE \), by
continuity.
The set \( \cE_{+} \) is the set of all nonnegative measurable functions.
 \end{proposition}

Now we extend the notion of integral to a wider class of functions.

% Gacs source line 3392
 \begin{definition}
  \label{def:g2.3.27}
  \notready
A measurable function \( f \) is called \df{integrable} with respect to a finite
measure \( \mu \) if \( \mu |f|^{+} < \infty \) and \( \mu |f|^{-} < \infty \).
In this case, we define \( \mu f = \mu |f|^{+} - \mu |f|^{-} \).
 \end{definition}

It is easy to see that the
mapping \( f\mapsto \mu f \) is linear when \( f \) runs through the set of
measureable functions with \( \mu |f|<\infty \).

The following is also easy to check.

% Gacs source line 3404
\begin{proposition}
  \label{prop:g2.3.28}
  \notready
Let \( \mu \) be a computable measure.
\begin{alphenum}
   \item
If \( f \) is computable then a program to compute
\( \mu f \) can be found from the program to compute \( f \).
   \item
If \( f \) is lower semicomputable then a program to lower semicompute
\( \mu f \) can be found from a program to lower semicompute \( f \).
\end{alphenum}
\end{proposition}

\subsection{Randomness tests}

We can now define randomness tests similarly to Section~\ref{sec:discrand}.

% Gacs source line 3441
 \begin{definition}
  \label{def:g2.3.29}
  \lean{Kolmogorov.Randomness.IsIntegrableSeqTest, Kolmogorov.Randomness.IsMLSeqTest, Kolmogorov.Randomness.IsUniversalIntegrableSeqTest, Kolmogorov.Randomness.seqVal, Kolmogorov.Randomness.levelIntegral}
  \leanok
  \uses{def:g2.3.5, def:g2.3.8, def:semicomputability}
A function \( d:X\to\dR \) is an \df{integrable test}, or \df{expectation-bounded test}
of randomness with respect to the probability measure \( P \)
if it is lower semicomputable and satisfies the condition
 \begin{equation} 	\label{eq:newtestcond-seq}
   \int 2^{d(x)}P(dx) \le 1.
 \end{equation}
It is called a \df{Martin-L\"of test}, or \df{probability-bounded test},
if instead of the latter condition
only the weaker one is satisfied saying \( P(d(x)>m)<2^{-m} \) for each
positive integer \( m \).

It is \df{universal} if it dominates all other integrable tests to within an
additive constant.
Universal Martin-L\"of tests are defined in the same way.
 \end{definition}

\begin{formalization}
A test on sequences is carried by a function \( T \) from strings to \( [0,\infty] \) that is
lower semicomputable and monotone along prefixes. Its value on \( \xi \) is
\( 2^{d(\xi)} = \sup_{n} T(\xi_{1:n}) \) (\texttt{seqVal T \( \xi \)}). The integral
in~\eqref{eq:newtestcond-seq} becomes the supremum of the level sums
\( \sum_{\len{z} = n} P(z) T(z) \) (\texttt{levelIntegral}), so \( T \) is integrable when
every level sum is at most \( 1 \). The Martin-L\"of condition is
\( P\setof{\xi : 2^{d(\xi)} > 2^{k}} \le 2^{-k} \) for every \( k \in \dN \), the set being the
open set generated by \( \setof{z : T(z) > 2^{k}} \) (Definition~\ref{def:g2.3.8}).
Universality is \( 2^{d} \le 2^{c} \cdot 2^{d_{0}} \) along every sequence. Universal
Martin-L\"of tests are not defined, since no result uses them.
\end{formalization}

Randomness tests and constructive null sets are closely related.

% Gacs source line 3460
\begin{theorem}
  \label{thm:g2.3.2}
  \lean{Kolmogorov.Randomness.isConstructiveNull_seqVal_top, Kolmogorov.Randomness.exists_integrableSeqTest_of_constructiveNull, Kolmogorov.Randomness.IsIntegrableSeqTest.isMLSeqTest}
  \leanok
  \uses{def:g2.3.29, def:g2.3.13}
  Let \( P \) be a computable probability measure
over the set \( S=\Sg^{\infty} \).
There is a correspondence between constructive null sets and
randomness tests:
  \begin{alphenum}
    \item
  For every randomness tests \( d \), the set \( N=\setof{\xi: d(\xi)=\infty} \) is a
  constructive null set.
    \item
For every constructive null set \( N \) there is a randomness test \( d \) with
\( N=\setof{\xi: d(\xi)=\infty} \).
  \end{alphenum}
\end{theorem}
\begin{proof}
  \proves{thm:g2.3.2}
  \leanok
Let \( d \) be a randomness test: then for each \( k \) the set \( G_{k}=\setof{\xi: d(\xi)>k} \) is a
constructive open set, and by Markov's inequality
(wich is proved in the continuous case just as in the discrete case)
 we have \( P(G_{k})\le 2^{-k} \).
The sets \( G_{k} \) witness that \( N \) is a constructive null set.

Let \( N \) be a constructive null set with \( N\sbsq \bigcap_{k=1}^{\infty} G_{k} \),
where \( G_{k} \) is a uniform sequence of constructive open sets with
\( P(G_{k})=2^{-k} \).
Without loss of generality assume that the sequence \( G_{k} \) is decreasing.
Then the function \( d(\xi)=\sup\setof{k: \xi\in G_{k}} \) is lower semicomputable
and satisfies \( P\setof{\xi:d(\xi)\ge k} \le 2^{-k} \), so it is a Martin-L\"of test.
Just as in Proposition~\ref{propo:ml-integrable},
it is easy to check that \( d(x)-2\log d(x) -c \) is an integrable test for some constant \( c \).
\end{proof}

\begin{formalization}
Part~(a) is proved for every Martin-L\"of test, hence for every integrable test by Markov's
inequality (\texttt{IsIntegrableSeqTest.isMLSeqTest}). Part~(b) is proved with \( \sbsq \) in
place of \( = \), which is the direction Theorem~\ref{thm:univ-seq-test} needs, and with an
integrable test built directly: if \( N \sbsq \bigcap_{m} G_{m} \) with
\( P(G_{m}) \le 2^{-m} \), then \( T(z) = \sum_{m} 2^{m} \cdot 1[z \in G_{2m+2}] \), read on
strings, has every level sum at most \( \sum_{m} 2^{m} 2^{-2m-2} < 1 \) and is infinite on
\( N \). The correction \( d - 2\log d - c \) is not needed.
\end{formalization}

Just as there is a universal constructive null set, there are universal
randomness tests.

% Gacs source line 3493
\begin{theorem}\label{thm:univ-integrable-test-seq}
  \lean{Kolmogorov.Randomness.canonicalSeqTest, Kolmogorov.Randomness.canonicalSeqTest_isUniversal, Kolmogorov.Randomness.seqVal_le_exact, Kolmogorov.Randomness.seqVal_le_canonicalSeqTest_exact, Kolmogorov.Randomness.seqVal_le_exact'}
  \leanok
  \uses{def:g2.3.29, def:ch2-description}
For a computable measure \( P \), there is a universal integrable test \( \ol\d_{P}(\xi) \).
More exactly, the function \( \xi\mapsto\ol\d_{P}(\xi) \) is lower semicomputable, satisfies~\eqref{eq:newtestcond} and
for all integrable tests \( d(\xi) \) for \( P \), we have
 \begin{equation*}d(\xi) \lea \ol\d_{P}(\xi) + K(d) + K(P).
 \end{equation*}
 \end{theorem}
The proof is similar to the proof Theorem~\ref{thm:univsem}
on the existence of a universal constructive semimeasure.
It uses the projection technique and a weighted sum.

\begin{proof}
  \proves{thm:univ-integrable-test-seq}
  \leanok
  \uses{thm:defic-charac-cpt-seqs, thm:ch2-exact-domination}
\end{proof}

\begin{formalization}
The universal test is explicit, not a projection over an enumeration of tests:
\( T_{P}(z) = \max_{y \prefix z} 2^{-K(y)}/P(y) \) (\texttt{canonicalSeqTest U P}), whose value
along \( \xi \) is \( \sup_{n} 2^{-K(\xi_{1:n})}/P(\xi_{1:n}) \). It is integrable by the first
half of the proof of Theorem~\ref{thm:defic-charac-cpt-seqs}, and it dominates every
integrable test by the second half. With his constant:
\( 2^{-K(y_{d})} 2^{-K(y_{P})} 2^{d(\xi)} \le 2^{c} \cdot 2^{\ol\d_{P}(\xi)} \) for
descriptions \( y_{d} \) of the test and \( y_{P} \) of \( P \), where \( c \) depends on
neither \( P \), \( d \) nor \( \xi \) (\texttt{seqVal\_le\_canonicalSeqTest\_exact}). Its proof
reruns the second half uniformly in \( \ang{y_{d}, y_{P}} \) and ends with
Theorem~\ref{thm:ch2-exact-domination}. The condition~\eqref{eq:newtestcond} in the statement
is to be read as~\eqref{eq:newtestcond-seq}.
\end{formalization}

\subsection{Randomness and complexity}

We hope for a result connecting complexity with randomness similar to
Theorem~\ref{thm:univ-integrable-test}.
Somewhat surprisingly, there is indeed such a result.
This is a surprise since in an infinite sequence, arbitrarily large oscillations
of any quantity depending on the prefixes are actually to be expected.

% Gacs source line 3513
 \begin{theorem}\label{thm:defic-charac-cpt-seqs}
  \lean{Kolmogorov.Randomness.universalSeqTest_equiv, Kolmogorov.Randomness.complexityRatioSup, Kolmogorov.Randomness.seqVal_le_of_isIntegrableSeqTest, Kolmogorov.Randomness.canonicalSeqTest_isIntegrable, Kolmogorov.Randomness.levelIntegral_canonicalSeqTest_le}
  \leanok
  \uses{def:g2.3.29, def:K}
Let \( X=\Sg^{\dN} \) be the set of infinite sequences.
For all computable measures \( \mu \) over \( X \), we have
 \begin{equation}\label{eq:defic-charac-seq}
 \ol\d_{\mu}(\xi) \eqa \sup_{n}\Paren{-\log \mu(\xi_{1:n}) - K(\xi_{1:n})}.
 \end{equation}
Here, the constant in \( \eqa \) depends on the computable measure \( \mu \).
 \end{theorem}

% Gacs source line 3522
 \begin{corollary}
  \label{cor:g2.3.30}
  \lean{Kolmogorov.Randomness.uniform_universalSeqTest_equiv, Kolmogorov.Randomness.complexityRatioSup_uniform, Kolmogorov.Randomness.isRandomSeq_uniform_iff}
  \leanok
  \uses{def:g2.3.29, def:g2.3.5}
Let \( \lg \) be the uniform distribution over the set of infinite binary sequences.
Then
 \begin{equation}
 \ol\d_{\lg}(\xi) \eqa \sup_{n}\Paren{n - K(\xi_{1:n})}.
 \end{equation}
 \end{corollary}
\begin{proof}
  \proves{cor:g2.3.30}
  \leanok
  \uses{thm:defic-charac-cpt-seqs, thm:univ-seq-test}
\end{proof}
In other words, an infinite binary sequence is random (with respect to the
uniform distribution) if and only if
the complexity \( K(\xi_{1:n}) \) of its initial
segments of length \( n \) never decreases below \( n \) by more than an
additive constant.

\begin{formalization}
For \( \lg \) the supremum is \( \sup_{n} 2^{n} \cdot 2^{-K(\xi_{1:n})} \)
(\texttt{complexityRatioSup\_uniform}). The statement in words is the Levin--Schnorr theorem,
\texttt{isRandomSeq\_uniform\_iff}: \( \xi \) is random for \( \lg \) if and only if there is a
\( c \) with \( n \le K(\xi_{1:n}) + c \) for all \( n \).
\end{formalization}

% Gacs source line 3534
 \begin{proof}[Proof of Theorem~\protect\ref{thm:defic-charac-cpt-seqs}]
  \proves{thm:defic-charac-cpt-seqs}
  \leanok
  \uses{thm:coding, l:kraft}
To prove \( \lea \), define the function
 \[
 f_{\mu}(\xi) = \sum_{s} 1_{sX}(\xi) \frac{\m(s \mvert \mu)}{\mu(s)}
  = \sum_{n} \frac{\m(\xi_{1:n} \mvert \mu)}{\mu(\xi_{1:n})}
  \ge \sup_{n} \frac{\m(\xi_{1:n} \mvert \mu)}{\mu(\xi_{1:n})}.
 \]
The function \( \xi \mapsto f_{\mu}(\xi) \)
is lower semicomputable with \( \mu^{\xi} f_{\mu}(\xi) \le 1 \), and hence
 \[
  \ol\d_{\mu}(\xi) \gea \log f(\xi) \gea
   \sup_{n}\Paren{-\log\mu(\xi_{1:n}) - K(\xi_{1:n} \mvert \mu)}.
 \]

The proof of \( \lea \) reproduces the proof of Theorem 5.2
of~\cite{GacsExact80}.
Let us abbreviate:
 \begin{equation*}1_{y}(\xi) = 1_{yX}(\xi).
 \end{equation*}
Since \( \cei{d(\xi)} \) only takes integer values and is lower semicomputable,
there are computable
sequences \( y_{i} \in \Sg^{*} \) and \( k_{i} \in \dN \) with
 \begin{align*}
   2^{\cei{d(\xi)}} &= \sup_{i}\; 2^{k_{i}} 1_{y_{i}}(\xi) \ge
  \frac{1}{2}\sum_{i} 2^{k_{i}} 1_{y_{i}}(\xi)
 \end{align*}
with the property that if \( i<j \) and \( 1_{y_{i}}(\xi)=1_{y_{j}}(\xi)=1 \) then
\( k_{i}<k_{j} \).
The inequality follows from the fact that for any finite sequence
\( n_{1}<n_{2}<\dots \), \( \sum_{j} 2^{n_{j}} \le 2 \max_{j} 2^{n_{j}} \).
The function \( \gm(y)=\sum_{y_{i}=y}2^{k_{i}} \) is lower semicomputable.
With it, we have
\begin{align*}
\sum_{i} 2^{k_{i}} 1_{y_{i}}(\xi)  &= \sum_{y\in\Sg^{*}}1_{y}(\xi)\gm(y).
 \end{align*}
Since \( \mu 2^{\cei{d}} \le 2 \), we have \( \sum_{y}\mu(y)\gm(y) \le 2 \), hence
\( \mu(y)\gm(y)\lem\m(y) \), that is \( \gm(y)\lem \m(y)/\mu(y) \).
It follows that
 \[
   2^{d(\xi)} \lem \sup_{y \in \Sg^{*}}\; 1_{y}(\xi)\frac{\m(y)}{\mu(y)}
  = \sup_{n} \frac{\m(\xi_{1:n})}{\mu(\xi_{1:n})}.
 \]
Taking logarithms, we obtain the \( \lea \) part of the theorem.
 \end{proof}

\begin{formalization}
Formalized multiplicatively, with unconditional \( K(\xi_{1:n}) \) in place of
\( K(\xi_{1:n} \mvert \mu) \): every universal integrable test \( T_{0} \) satisfies
\( 2^{-c} R(\xi) \le \sup_{n} T_{0}(\xi_{1:n}) \le 2^{c} R(\xi) \), where
\( R(\xi) = \sup_{n} 2^{-K(\xi_{1:n})}/\mu(\xi_{1:n}) \) and \( c \) depends on \( \mu \)
(\texttt{universalSeqTest\_equiv}).
\begin{itemize}
\item The first half (his first ``\( \lea \)'' should read ``\( \gea \)'') is that
  \( T_{\mu}(z) = \max_{y \prefix z} 2^{-K(y)}/\mu(y) \) is an integrable test
  (\texttt{canonicalSeqTest\_isIntegrable}). A level sum regroups by prefixes into
  \( \sum_{y} 2^{-K(y)} \le 1 \).
\item The second half, that every integrable test \( T \) has
  \( \sup_{n} T(\xi_{1:n}) \le 2^{c} R(\xi) \)
  (\texttt{seqVal\_le\_of\_isIntegrableSeqTest}), follows his decomposition with the
  enumeration made explicit. Level \( k \) is \( \setof{T > 2^{k+1}} \). A string \( z \)
  hits level \( k \) if the stage-\( \len{z} \) approximation of \( T \), maximized over the
  prefixes of \( z \), exceeds \( 2^{k+1} \); it is a first hit if no proper prefix hits. For
  each \( k \) the first hits form a decidable prefix-free set, so
  \( \gm(z) = \sum\setof{2^{k} : z \text{ is a first hit of level } k} \) is computable. A
  layer-cake bound gives \( \sum_{z} \mu(z) \gm(z) \le 1 \), and the coding theorem gives
  \( \gm \le 2^{c} \cdot 2^{-K}/\mu \).
\end{itemize}
\end{formalization}

\subsection{Universal semimeasure, algorithmic probability}

Universal semimeasures exist over infinite sequences just as in the discrete
space.

% Gacs source line 3584
\begin{definition}
  \label{def:g2.3.31}
  \lean{Kolmogorov.Randomness.IsContSemimeasure}
  \leanok
  \uses{def:semicomputability}
  A function \( \mu:\Sg^{*}\to\cR_{+} \) is a \df{semimeasure} if it satisfies
 \begin{align*}
     \mu(x)      &\ge\sum_{s\in\Sg}\mu(xs),
\\ \mu(\Lg) &\le 1.
 \end{align*}
If it is lower semicomputable we will call the semimeasure \df{constructive}.
\end{definition}
These requirements imply, just as for measures,
the following generalization to an arbitrary prefix-free
set \( A\sbsq \Sg^{*} \):
 \begin{align*}
   \sum_{x\in A}\mu(x)\le 1.
 \end{align*}

\begin{formalization}
For \( \Sg = \dB \), \texttt{IsContSemimeasure \( \mu \)} is \( \mu(\Lg) \le 1 \) and
\( \mu(x0) + \mu(x1) \le \mu(x) \); constructive means, in addition, lower semicomputable.
The bound for prefix-free sets is formalized for finite \( A \)
(\texttt{sum\_le\_one\_of\_prefixFree}); the infinite case follows by taking suprema but is
not stated.
\end{formalization}

Standard technique gives the following:

% Gacs source line 3601
\begin{theorem}[Universal semimeasure]
  \label{thm:g2.3.5}
  \lean{Kolmogorov.Randomness.univContSemimeasure_isUniversal, Kolmogorov.Randomness.exists_universalContSemimeasure, Kolmogorov.Randomness.IsUniversalContSemimeasure, Kolmogorov.Randomness.contSemimeasure_le_univ_exact}
  \leanok
  \uses{def:g2.3.31}
  There is a universal constructive semimeasure \( \mu \), that is a
constructive semimeasure with the property that for every other constructive
semimeasure \( \nu \) there is a constant \( c_{\nu}>0 \) with
\( \mu(x)\ge c_{\nu}\nu(x) \).
\end{theorem}

\begin{proof}
  \proves{thm:g2.3.5}
  \leanok
  \uses{thm:univsem, thm:ch2-exact-domination}
\end{proof}

\begin{formalization}
The construction is that of Theorem~\ref{thm:univsem}:
\( M(x) = \sum_{i} 2^{-i-1} \nu_{i}(x) \), where \( \nu_{i} \) is the \( i \)-th lower
semicomputable function, trimmed into a continuous semimeasure that equals it whenever it
already is one. His trimming, which freezes an approximation once it violates the
constraint, does not carry over: the lower approximations of a continuous semimeasure can
violate \( \mu(x0) + \mu(x1) \le \mu(x) \) at finite stages. The formalization uses the
construction of Zvonkin and Levin~\cite{ZvLe70} instead. At stage \( s \) it raises the
values bottom-up, \( \varphi(x) = \max(\psi_{s}(x), \varphi(x0) + \varphi(x1)) \) on the
strings of length at most \( s \), and freezes only when the root exceeds \( 1 \)
(\texttt{trimSt}). The exact form \( 2^{-K(y)} \nu(x) \le 2^{c} M(x) \) for a description
\( y \) of \( \nu \), with \( c \) independent of \( \nu \), is
\texttt{contSemimeasure\_le\_univ\_exact}.
\end{formalization}

% Gacs source line 3608
\begin{definition}
  \label{def:g2.3.32}
  \lean{Kolmogorov.Randomness.univContSemimeasure, Kolmogorov.Randomness.IsUniversalContSemimeasure}
  \leanok
  \uses{thm:g2.3.5}
  Let us fix a universal constructive semimeasure over \( X \) and denote it by
\( M(x) \).
\end{definition}

\begin{formalization}
No semimeasure is fixed: as with \( U \) in Chapter~1, results take any \( M \) with
\texttt{IsUniversalContSemimeasure M} as a hypothesis, and \texttt{univContSemimeasure} is one.
\end{formalization}

There is a graphical way to represent a universal semimeasure, via monotonic
Turing machines, which can be viewed a generalization of self-delimiting
machines.

% Gacs source line 3617
\begin{definition}
  \label{def:g2.3.33}
A Turing machine \( \cT \) is \df{monotonic} if it
has no input tape or output tape, but can ask repeatedly for input symbols and can emit
repeatedly output symbols.
The input is assumed to be an infinite binary string, but it is not assumed that
\( \cT \) will ask for all of its symbols, even in infinite time.
If the finite or
infinite input string is \( p=p_{1}p_{2}\dotsm \), the (finite or infinite) output string is
written as \( T(p)\in\Sg^{*}\cup\Sg^{\dN} \).
\end{definition}

\begin{formalization}
Not formalized yet. Monotonic machines, Definitions~\ref{def:g2.3.34} and~\ref{def:g2.3.35}
and Theorem~\ref{thm:g2.3.6} are planned future work. Nothing else in the chapter depends on
them.
\end{formalization}

Just as we can generalize self-delimiting machines to obtain the notion of
monotonic machines, we can generalize algorithmic probability, defined for a
discrete space, to a version defined for infinite sequences.
Imagine the monotonic Turing machine \( \cT \) computing the function \( T(p) \)
as in the definition, and assume that its input symbols come from coin-tossing.

% Gacs source line 3634
\begin{definition}[Algorithmic probability over infinite sequences]
  \label{def:g2.3.34}
  \uses{def:g2.3.33}
For a string \( x\in\Sg^{*} \), let \( P_{T}(x) \) be the probability of \( T(\pi)\postfix x \),
when \( \pi=\pi_{1}\pi_{2}\dotsm \) is the infinite coin-tossing sequence.
\end{definition}

We can compute
 \begin{align*}
  P_{T}(x)=\sum_{T(p) \postfix x, p\text{ minimal}} 2^{-\len{p}},
 \end{align*}
where ``minimal'' means that no prefix \( p' \) of \( p \) gives \( T(p')\postfix x \).
This expression shows that \( P_{T}(x) \) is lower semicomputable.
It is also easy to check that it is a semimeasure, so we have a constructive
semimeasure.
The following theorem is obtained using a standard construction:

% Gacs source line 3649
\begin{theorem}
  \label{thm:g2.3.6}
  \uses{def:g2.3.34, def:g2.3.31}
Every constructive semimeasure \( \mu(x) \) can be represented as
\( \mu(x)=P_{T}(x) \) with the help of an appropriate a monotonic machine \( \cT \).
\end{theorem}

\begin{formalization}
Not formalized yet; planned future work. This is the one hard construction of the chapter, a
monotone analogue of the Kraft--Chaitin theorem. No other result depends on it: the
formalization works with \( M \) as a universal constructive semimeasure
(Theorem~\ref{thm:g2.3.5}), never with a machine.
\end{formalization}

This theorem justifies the following.

% Gacs source line 3656
\begin{definition}
  \label{def:g2.3.35}
  \uses{thm:g2.3.6, def:g2.3.32}
From now on we will call the universal semimeasure \( M(x) \) also
the \df{algorithmic probability}.
A monotonic Turing machine \( \cT \) giving \( P_{T}(x)=M(x) \) will be called an
\df{optimal machine}.
\end{definition}

\begin{formalization}
Not formalized yet: an optimal machine needs Theorem~\ref{thm:g2.3.6}. In the formalization
the algorithmic probability is \( M \) of Definition~\ref{def:g2.3.32}.
\end{formalization}

We will see that \( -\log M(x) \) should also be considered a kind of
description complexity:

% Gacs source line 3666
  \begin{definition}
  \label{def:g2.3.36}
  \lean{Kolmogorov.Randomness.univContSemimeasure}
  \leanok
  \uses{def:g2.3.32}
Denote
 \begin{align*}
   KM(x) =-\log M(x).
 \end{align*}
  \end{definition}

\begin{formalization}
No \( KM \) is defined in Lean. Every statement about \( KM \) is made about
\( M = 2^{-KM} \).
\end{formalization}

How does \( KM(x) \) compare to \( K(x) \)?

% Gacs source line 3675
\begin{proposition}\label{prop:KmHCompare}
  \lean{Kolmogorov.Randomness.KM_bounds, Kolmogorov.Randomness.prefixComplexityWeight_le_univ, Kolmogorov.Randomness.univ_mul_length_le_prefixComplexityWeight, Kolmogorov.Randomness.extSum}
  \leanok
  \uses{def:g2.3.36, def:K}
We have the bounds
 \begin{align*}
 KM(x) \lea K(x) \lea KM(x) + K(\len{x}).
 \end{align*}
\end{proposition}
\begin{proof}
  \proves{prop:KmHCompare}
  \leanok
  \uses{thm:g2.3.5, thm:coding, l:kraft}
To prove \( KM(x)\lea K(x) \) define the lower semicomputable function
 \begin{align*}
   \mu(x) = \sup_{S}\sum_{y\in S}\m(xy),
 \end{align*}
By its definition this is a constructive semimeasure and it is
\( \ge \m(x) \).
This implies \( M(x)\gem \m(x) \), and thus \( KM(x)\lea K(x) \).

For the other direction, define the following lower semicomputable function \( \nu \)
over \( \Sg^{*} \):
 \begin{align*}
   \nu(x)=\m(\len{x})\cdot M(x).
 \end{align*}
To show that it is a semimeasure compute:
 \begin{align*}
   \sum_{x\in\Sg^{*}}\nu(x)=\sum_{n}\m(n)\sum_{x\in\Sg^{n}}M(x)\le\sum_{n}\m(n)\le 1.
 \end{align*}
It follows that \( \m(x)\gem\nu(x) \) and hence \( K(x) \lea KM(x)+K(\len{x}) \).
\end{proof}

\begin{formalization}
Multiplicatively, with one constant: \( 2^{-K(x)} \le 2^{c} M(x) \) and
\( M(x) \cdot 2^{-K(\len{x})} \le 2^{c} \cdot 2^{-K(x)} \) (\texttt{KM\_bounds}). For the first
bound the formalization takes \( \mu(x) = \sum_{y} 2^{-K(xy)} \), the sum over all
extensions of \( x \) (\texttt{extSum}), instead of a supremum over sets \( S \). It is a
constructive continuous semimeasure because \( \sum_{x} 2^{-K(x)} \le 1 \).
\end{formalization}

\subsection{Randomness via algorithmic probability}

Fix a computable measure \( P(x) \) over the space \( X \).
Just as in the discrete case, the universal semimeasure gives rise to a
Martin-L\"of test.

% Gacs source line 3710
\begin{definition}
  \label{def:g2.3.38}
  \lean{Kolmogorov.Randomness.ratioTest, Kolmogorov.Randomness.seqVal_ratioTest}
  \leanok
  \uses{def:g2.3.32, def:g2.3.5, def:g2.3.29}
Denote
 \begin{align*}
 \d'_{P}(\xi) = \log\sup_{n}\frac{M(\xi_{1:n})}{P(\xi_{1:n})}.
 \end{align*}
\end{definition}

By Proposition~\ref{prop:KmHCompare}, this function is
is \( \gem\ol\d_{P}(\xi) \), defined in Theorem~\ref{thm:univ-integrable-test-seq}.
It can be shown that it is no longer expectation-bounded.
But the theorem below shows that it is still a Martin-L\"of
(probability-bounded) test, so it defines the same infinite random sequences.

\begin{formalization}
\texttt{ratioTest \( \nu \) P} is defined for any continuous semimeasure \( \nu \) in place of
\( M \), as the prefix maximum \( z \mapsto \max_{y \prefix z} \nu(y)/P(y) \); for
\( \nu = M \) its value along \( \xi \) is \( 2^{\d'_{P}(\xi)} \). The comparison with
\( \ol\d_{P} \) is \texttt{complexityRatioSup\_le\_ratioTest}, from the first bound of
Proposition~\ref{prop:KmHCompare}. That \( \d'_{P} \) is in general not expectation-bounded is
not formalized.
\end{formalization}

% Gacs source line 3723
\begin{theorem}
  \label{thm:g2.3.7}
  \lean{Kolmogorov.Randomness.ratioTest_isMLSeqTest, Kolmogorov.Randomness.openMeasure_ratioTest_le, Kolmogorov.Randomness.complexityRatioSup_le_ratioTest, Kolmogorov.Randomness.isRandomSeq_iff_ratioTest}
  \leanok
  \uses{def:g2.3.38, def:g2.3.29}
The function \( \d'_{P}(\xi) \) is a Martin-L\"of test.
\end{theorem}
\begin{proof}
  \proves{thm:g2.3.7}
  \leanok
  \uses{def:g2.3.31, cor:g2.3.7}
Lower semicomputability follows from the form of definition.
By standard methods,
produce a list of finite strings \( y_{1},y_{2},\dots \) with:
\begin{alphenum}
\item \( M(y_{i})/P(y_{i})>2^{m} \)
\item The cylinder sets \( y_{i}X \) cover of \( G_{m} \), and are disjoint (thus the
  set \( \{y_{1},y_{2},\dots\} \) is prefix-free).
\end{alphenum}
We have
 \begin{align*}
   P(G_{m})=\sum_{i}P(y_{i})<2^{-m}\sum_{i}M(y_{i})\le 2^{-m}.
 \end{align*}
\end{proof}

\begin{formalization}
The ``standard methods'' are made explicit. For each \( n \), the prefix-minimal strings
\( y \) of length at most \( n \) with \( M(y)/P(y) > 2^{m} \) form a finite prefix-free set,
so their total \( P \)-measure is at most \( 2^{-m} \sum_{i} M(y_{i}) \le 2^{-m} \)
(\texttt{openMeasure\_ratioTest\_le}). The bound holds for every continuous semimeasure in
place of \( M \); lower semicomputability is used only to make \( \d'_{P} \) lower
semicomputable. That \( \d'_{P} \) defines the same random sequences as \( \ol\d_{P} \) is
\texttt{isRandomSeq\_iff\_ratioTest}.
\end{formalization}

Just as in the discrete case, the universality of \( M(x) \) implies
 \begin{align*}
   KM(x) \lea -\log P(x) +K(P).
 \end{align*}
      On the other hand,
 \begin{align*}
   \sup_{n} -\log P(\xi_{1:n})-KM(\xi_{1:n})
 \end{align*}
is a randomness test.
So for random \( \xi \), the value \( KM(\xi_{1:n}) \) remains within a constant of
\( -\log P(x) \): it does not oscillate much.

\begin{formalization}
\( KM(x) \lea -\log P(x) \), with a constant depending on \( P \), is
\texttt{seqMeasure\_le\_univ}. The form with \( K(P) \) follows from
\texttt{contSemimeasure\_le\_univ\_exact} (Theorem~\ref{thm:g2.3.5}) applied to a description
of \( P \), but is not stated. That \( KM(\xi_{1:n}) \) stays within a constant of
\( -\log P(\xi_{1:n}) \) along a random \( \xi \) is \texttt{univ\_comparable\_on\_random}.
\end{formalization}
